\documentclass{beamer}
\usetheme{Madrid}
\usecolortheme{default}
\usepackage[utf8]{inputenc}
\usepackage{amsmath}
\usepackage{amsfonts}
\usepackage{amssymb}
\usepackage{tikz}

\title{CRISC Hard Questions from ISACA QAE Manual}
\subtitle{200 Selected Questions --- Answers on the Next Page}
\author{Compiled from ISACA CRISC Review Questions, Answers \& Explanations Manual, 6th Edition}
\date{\today}

\begin{document}
	
	\begin{frame}
		\titlepage
	\end{frame}
	
	\begin{frame}{Outline}
		\tableofcontents
	\end{frame}
	
	% ============================================================
	% DOMAIN 1 --- GOVERNANCE (26%) --- Q1--Q55
	% ============================================================
	\section{Domain 1: Governance}
	
	% ---- D1 Q1 ----
	\begin{frame}{D1-Q1: Defining Risk Management Strategies}
		\begin{block}{Question}
			Which of the following is MOST important to determine when defining risk management strategies?
		\end{block}
		\begin{enumerate}[(A)]
			\item Risk assessment criteria
			\item IT architecture complexity
			\item Enterprise disaster recovery plan
			\item Business objectives and operations
		\end{enumerate}
	\end{frame}
	\begin{frame}{D1-Q1: Answer}
		\begin{block}{Correct Answer: D --- Business objectives and operations}
			The risk practitioner must analyze the enterprise's objectives and risk tolerance and define a risk management framework based on this analysis. Some enterprises accept known risk; others invest in mitigating controls. Risk assessment criteria alone are insufficient; IT complexity relates to assessment, not strategy; a DR plan relates to mitigation.
		\end{block}
	\end{frame}
	
	% ---- D1 Q2 ----
	\begin{frame}{D1-Q2: Risk Management Strategic Plan}
		\begin{block}{Question}
			Which of the following is the MOST important information to include in a risk management strategic plan?
		\end{block}
		\begin{enumerate}[(A)]
			\item Risk management staffing requirements
			\item Risk management mission statement
			\item Risk mitigation investment plans
			\item Current state and desired future state
		\end{enumerate}
	\end{frame}
	\begin{frame}{D1-Q2: Answer}
		\begin{block}{Correct Answer: D --- Current state and desired future state}
			It is most important to paint a vision for the future and draw a road map from the starting point, which requires that the current state and desired future state be fully understood. Staffing, mission, and investment plans are generally driven by that understanding.
		\end{block}
	\end{frame}
	
	% ---- D1 Q3 ----
	\begin{frame}{D1-Q3: Business Unit Risk Roles}
		\begin{block}{Question}
			Which BEST describes the risk-related roles of an organizational business unit (BU)? The BU management team:
		\end{block}
		\begin{enumerate}[(A)]
			\item Owns the mitigation plan while board members identify and assess risk
			\item Owns the risk and is responsible for identifying, assessing, mitigating, and reporting
			\item Carries out risk responsibilities, but ultimate accountability belongs to the board
			\item Is ultimately accountable, and the board owns the risk
		\end{enumerate}
	\end{frame}
	\begin{frame}{D1-Q3: Answer}
		\begin{block}{Correct Answer: B}
			The BU is responsible for owning the risk and its resulting actions. Risk owners identify, measure, monitor, control, and report on risk to executive management as established by the corporate risk framework. The board does not perform risk identification or assessment functions.
		\end{block}
	\end{frame}
	
	% ---- D1 Q4 ----
	\begin{frame}{D1-Q4: Outsourced IT in Foreign Country}
		\begin{block}{Question}
			An enterprise has outsourced the majority of its IT department to a third party whose servers are in a foreign country. Which is the MOST critical security consideration?
		\end{block}
		\begin{enumerate}[(A)]
			\item Security breach notification may be delayed due to the time difference
			\item Additional network intrusion detection sensors are required
			\item The enterprise could be unable to monitor compliance with internal security and privacy guidelines
			\item Laws and regulations of the country of origin may not be enforceable in the foreign country
		\end{enumerate}
	\end{frame}
	\begin{frame}{D1-Q4: Answer}
		\begin{block}{Correct Answer: D}
			Laws and regulations of the country of origin may not be enforceable in the foreign country. Conversely, the laws of the foreign vendor may affect the enterprise. Potential violation of local laws may not be recognized or remedied due to lack of knowledge and/or inability to enforce.
		\end{block}
	\end{frame}
	
	% ---- D1 Q5 ----
	\begin{frame}{D1-Q5: Improving Risk Management Process}
		\begin{block}{Question}
			Which of the following is MOST beneficial to improvement of an enterprise's risk management process?
		\end{block}
		\begin{enumerate}[(A)]
			\item Key risk indicators
			\item External benchmarking
			\item Latest risk assessment
			\item A maturity model
		\end{enumerate}
	\end{frame}
	\begin{frame}{D1-Q5: Answer}
		\begin{block}{Correct Answer: D --- A maturity model}
			A maturity model helps identify the status quo and the path to the desired state and is most helpful when an enterprise desires to improve a business process such as risk management. KRIs and benchmarking do not define the desired state; the latest risk assessment is only an input.
		\end{block}
	\end{frame}
	
	% ---- D1 Q6 ----
	\begin{frame}{D1-Q6: Independent Risk Review}
		\begin{block}{Question}
			Which is the PRIMARY reason for subjecting the risk management process to review by independent risk auditors and assessors?
		\end{block}
		\begin{enumerate}[(A)]
			\item To ensure that the risk results are consistent
			\item To ensure that the risk factors and risk profile are well-defined
			\item To correct any mistakes in risk assessment
			\item To validate the control weaknesses for management reporting
		\end{enumerate}
	\end{frame}
	\begin{frame}{D1-Q6: Answer}
		\begin{block}{Correct Answer: B}
			Risk profile and risk factors are defined during the risk assessment process; an independent review helps ensure that the underlying process is effective and helps identify areas for future improvement.
		\end{block}
	\end{frame}
	
	% ---- D1 Q7 ----
	\begin{frame}{D1-Q7: Evaluating Legal/Regulatory Impact}
		\begin{block}{Question}
			What is the MOST effective method to evaluate the potential impact of legal, regulatory, and contractual requirements on business objectives?
		\end{block}
		\begin{enumerate}[(A)]
			\item A compliance-oriented gap analysis
			\item Interviews with business process stakeholders
			\item A mapping of compliance requirements to policies and procedures
			\item A compliance-oriented business impact analysis
		\end{enumerate}
	\end{frame}
	\begin{frame}{D1-Q7: Answer}
		\begin{block}{Correct Answer: D}
			A compliance-oriented BIA identifies compliance requirements to which the enterprise is subject and assesses their effect on business objectives and activities. Gap analysis only finds gaps; interviews identify objectives but not compliance; mapping shows how compliance is achieved but not business impact.
		\end{block}
	\end{frame}
	
	% ---- D1 Q8 ----
	\begin{frame}{D1-Q8: Maturity Assessment RFP}
		\begin{block}{Question}
			An enterprise is hiring a consultant to determine the maturity level of the risk management program. The MOST important element of the RFP is the:
		\end{block}
		\begin{enumerate}[(A)]
			\item Sample deliverable
			\item Past experience of the engagement team
			\item Methodology used in the assessment
			\item References from other enterprises
		\end{enumerate}
	\end{frame}
	\begin{frame}{D1-Q8: Answer}
		\begin{block}{Correct Answer: C --- Methodology}
			Methodology illustrates the consultant's process and offers a basis to align expectations with execution of the assessment. It establishes requirements of all parties involved.
		\end{block}
	\end{frame}
	
	% ---- D1 Q9 ----
	\begin{frame}{D1-Q9: New Law Affecting HR System}
		\begin{block}{Question}
			Shortly after the annual policy review, a risk practitioner becomes aware that a new law may affect security requirements for the HR system. The risk practitioner should:
		\end{block}
		\begin{enumerate}[(A)]
			\item Analyze what systems and technology-related processes may be impacted
			\item Ensure necessary adjustments are implemented during the next review cycle
			\item Initiate an ad hoc revision of the corporate policy
			\item Notify the system custodian to implement changes
		\end{enumerate}
	\end{frame}
	\begin{frame}{D1-Q9: Answer}
		\begin{block}{Correct Answer: A}
			Assessing what systems and technology-related processes may be impacted is the best course of action. The analysis must also determine whether existing controls already address the new requirements. Policy changes require proper approval; waiting can create compliance risk.
		\end{block}
	\end{frame}
	
	% ---- D1 Q10 ----
	\begin{frame}{D1-Q10: Outsourced Data Processing Violation}
		\begin{block}{Question}
			An enterprise outsourced personnel data processing to a supplier, and a regulatory violation occurs during processing. Who will be held legally responsible?
		\end{block}
		\begin{enumerate}[(A)]
			\item The supplier, because it has operational responsibility
			\item The enterprise, because it owns the data
			\item The enterprise and the supplier
			\item The supplier, because it did not comply with the contract
		\end{enumerate}
	\end{frame}
	\begin{frame}{D1-Q10: Answer}
		\begin{block}{Correct Answer: B}
			The enterprise retains responsibility for management of and adherence to policies, procedures, and regulatory requirements. The regulatory authority will generally hold the enterprise responsible for failure to comply, including penalties. Supplier responsibility is limited to contract terms.
		\end{block}
	\end{frame}
	
	% ---- D1 Q11 ----
	\begin{frame}{D1-Q11: Evaluating Internal Controls Effectiveness}
		\begin{block}{Question}
			Which role is responsible for evaluating the effectiveness of existing internal information security controls within an enterprise?
		\end{block}
		\begin{enumerate}[(A)]
			\item Data owner
			\item Senior management
			\item End users
			\item System auditor
		\end{enumerate}
	\end{frame}
	\begin{frame}{D1-Q11: Answer}
		\begin{block}{Correct Answer: D --- System auditor}
			The system auditor provides continuous feedback to senior management about the effectiveness of internal controls. Data owners define requirements; senior management ensures effectiveness; end users are not responsible for evaluation.
		\end{block}
	\end{frame}
	
	% ---- D1 Q12 ----
	\begin{frame}{D1-Q12: Overall Effectiveness of Risk Program}
		\begin{block}{Question}
			Which of the following BEST ensures the overall effectiveness of a risk management program?
		\end{block}
		\begin{enumerate}[(A)]
			\item Obtaining feedback from all end users
			\item Assigning a dedicated risk manager to run the program
			\item Applying quantitative risk methodologies
			\item Gaining the participation of relevant stakeholders
		\end{enumerate}
	\end{frame}
	\begin{frame}{D1-Q12: Answer}
		\begin{block}{Correct Answer: D}
			Without participation from enterprise stakeholders and their active role in supervision and risk monitoring, the effectiveness of a risk management program cannot be guaranteed. All-user feedback is infeasible; a dedicated manager without stakeholders is less effective; methodology choice depends on needs.
		\end{block}
	\end{frame}
	
	% ---- D1 Q13 ----
	\begin{frame}{D1-Q13: Achieving Risk-Based Objectives}
		\begin{block}{Question}
			Which approach BEST helps an enterprise achieve risk-based organizational objectives?
		\end{block}
		\begin{enumerate}[(A)]
			\item Ensure that asset owners perform annual risk assessments
			\item Review and update the risk register regularly
			\item Assign a steering committee to the risk management process
			\item Embed risk management activities into business processes
		\end{enumerate}
	\end{frame}
	\begin{frame}{D1-Q13: Answer}
		\begin{block}{Correct Answer: D}
			The primary objective of embedding risk management activities into business processes is to achieve risk-based organizational objectives in the most effective manner possible. Assessments alone, register maintenance, or a steering committee are less effective.
		\end{block}
	\end{frame}
	
	% ---- D1 Q14 ----
	\begin{frame}{D1-Q14: Regulatory Reliance on Reviews}
		\begin{block}{Question}
			When assessing the capability of the risk management process, a regulatory body would place the GREATEST reliance on:
		\end{block}
		\begin{enumerate}[(A)]
			\item A peer review
			\item An internal review
			\item An external review
			\item A process capability review
		\end{enumerate}
	\end{frame}
	\begin{frame}{D1-Q14: Answer}
		\begin{block}{Correct Answer: C --- External review}
			Regulatory entities generally use assessments performed by an objective and independent third party. An external review is the most objective and independent of the choices presented.
		\end{block}
	\end{frame}
	
	% ---- D1 Q15 ----
	\begin{frame}{D1-Q15: Evaluating Current Risk Profile}
		\begin{block}{Question}
			A company has been improving its security and compliance program since the last review one year ago. What should it do to evaluate its current risk profile?
		\end{block}
		\begin{enumerate}[(A)]
			\item Review previous findings and ensure all issues are resolved
			\item Conduct follow-up audits in previously deficient areas
			\item Monitor KRIs and use results to develop targeted assessments
			\item Perform a new enterprise risk assessment using an independent expert
		\end{enumerate}
	\end{frame}
	\begin{frame}{D1-Q15: Answer}
		\begin{block}{Correct Answer: D}
			The best way to ensure the security posture is still compliant is to conduct another risk assessment. A lot can change in a year, and an independent expert provides more objective results than internal staff testing their own work.
		\end{block}
	\end{frame}
	
	% ---- D1 Q16 ----
	\begin{frame}{D1-Q16: Inadequate Data/System Ownership Policy}
		\begin{block}{Question}
			Which is the GREATEST risk of a policy that defines data and system ownership inadequately?
		\end{block}
		\begin{enumerate}[(A)]
			\item Audit recommendations may not be implemented
			\item Users may have unauthorized access to create, modify, or delete data
			\item User management coordination does not exist
			\item Specific user accountability cannot be established
		\end{enumerate}
	\end{frame}
	\begin{frame}{D1-Q16: Answer}
		\begin{block}{Correct Answer: B}
			Without a policy defining who grants access to specific data or systems, the risk that employees receive access without justified business purpose increases. Authority to grant access should be assigned to specific individuals.
		\end{block}
	\end{frame}
	
	% ---- D1 Q17 ----
	\begin{frame}{D1-Q17: Management's Role in Risk Strategy}
		\begin{block}{Question}
			Which BEST describes the role of management in implementing a risk management strategy?
		\end{block}
		\begin{enumerate}[(A)]
			\item Ensure that planning, budgeting, and performance of security components are appropriate
			\item Assess and incorporate the results of risk management into decision-making
			\item Identify, evaluate, and minimize risk to IT systems supporting the mission
			\item Understand the risk management process to develop training materials
		\end{enumerate}
	\end{frame}
	\begin{frame}{D1-Q17: Answer}
		\begin{block}{Correct Answer: B}
			Assessing and incorporating the results of risk management into the decision-making process best describes the role of senior management in establishing and implementing a risk management strategy.
		\end{block}
	\end{frame}
	
	% ---- D1 Q18 ----
	\begin{frame}{D1-Q18: Primary Focus of IT Business Risk}
		\begin{block}{Question}
			The PRIMARY focus of managing IT-related business risk is to protect:
		\end{block}
		\begin{enumerate}[(A)]
			\item Information
			\item Hardware
			\item Applications
			\item Databases
		\end{enumerate}
	\end{frame}
	\begin{frame}{D1-Q18: Answer}
		\begin{block}{Correct Answer: A --- Information}
			The primary objective for any enterprise is to protect mission-critical information based on a risk assessment. Hardware, applications, and databases are focuses only if they store, process, or transfer mission-critical data.
		\end{block}
	\end{frame}
	
	% ---- D1 Q19 ----
	\begin{frame}{D1-Q19: Best Risk Management Perspective}
		\begin{block}{Question}
			Which can provide the BEST perspective of risk management to an enterprise's employees and stockholders?
		\end{block}
		\begin{enumerate}[(A)]
			\item An interdisciplinary team within the enterprise
			\item A third-party risk assessment service provider
			\item The enterprise's IT department
			\item The enterprise's internal compliance department
		\end{enumerate}
	\end{frame}
	\begin{frame}{D1-Q19: Answer}
		\begin{block}{Correct Answer: A}
			An interdisciplinary team ensures all areas are adequately considered in risk assessment and provides an enterprise-wide perspective. Third parties lack internal knowledge; IT alone does not reflect the enterprise; compliance focuses on regulatory implications.
		\end{block}
	\end{frame}
	
	% ---- D1 Q20 ----
	\begin{frame}{D1-Q20: Global Policy for Multiple Regions}
		\begin{block}{Question}
			Which approach to corporate policy BEST supports an enterprise's expansion to other regions where different local laws apply?
		\end{block}
		\begin{enumerate}[(A)]
			\item A global policy without provisions that might be disputed locally
			\item A global policy amended to comply with local laws
			\item A global policy that complies with laws at enterprise headquarters
			\item Local policies to accommodate laws within each region
		\end{enumerate}
	\end{frame}
	\begin{frame}{D1-Q20: Answer}
		\begin{block}{Correct Answer: B}
			A global policy including local amendments ensures alignment with local laws and regulations. A single global policy cannot address all locales; HQ-only policy exposes legal risk; decentralized local policies are expensive and fail to leverage common practices.
		\end{block}
	\end{frame}
	
	% ---- D1 Q21 ----
	\begin{frame}{D1-Q21: Data Classification Responsibility}
		\begin{block}{Question}
			Who is MOST likely responsible for data classification?
		\end{block}
		\begin{enumerate}[(A)]
			\item Data user
			\item Data owner
			\item Data custodian
			\item System administrator
		\end{enumerate}
	\end{frame}
	\begin{frame}{D1-Q21: Answer}
		\begin{block}{Correct Answer: B --- Data owner}
			The data owner classifies data according to the enterprise's data classification scheme. The classification defines who is eligible to access the data and what controls are required. Custodians implement custody and storage.
		\end{block}
	\end{frame}
	
	% ---- D1 Q22 ----
	\begin{frame}{D1-Q22: Security Governance Model Impact}
		\begin{block}{Question}
			Which factor will have the GREATEST impact on the type of information security governance model an enterprise adopts?
		\end{block}
		\begin{enumerate}[(A)]
			\item The number of employees
			\item The enterprise's budget
			\item The organizational structure
			\item The type of technology used
		\end{enumerate}
	\end{frame}
	\begin{frame}{D1-Q22: Answer}
		\begin{block}{Correct Answer: C --- Organizational structure}
			Information security governance models depend significantly on the overall organizational structure. Employee count, budget, and technology do not primarily drive the governance model; well-defined processes provide proper governance.
		\end{block}
	\end{frame}
	
	% ---- D1 Q23 ----
	\begin{frame}{D1-Q23: Information Security Procedures}
		\begin{block}{Question}
			Information security procedures should:
		\end{block}
		\begin{enumerate}[(A)]
			\item Be updated frequently as new software is released
			\item Underline the importance of security governance
			\item Define the allowable limits of behavior
			\item Describe security baselines for each platform
		\end{enumerate}
	\end{frame}
	\begin{frame}{D1-Q23: Answer}
		\begin{block}{Correct Answer: A}
			Security procedures must change frequently to keep up with software changes. Because a procedure is a how-to document, it must be kept current. Governance is addressed in policy; behavioral limits are set by policy; baselines are set by standards.
		\end{block}
	\end{frame}
	
	% ---- D1 Q24 ----
	\begin{frame}{D1-Q24: Greatest Benefit of Risk-Aware Culture}
		\begin{block}{Question}
			Which is the GREATEST benefit of a risk-aware culture?
		\end{block}
		\begin{enumerate}[(A)]
			\item Issues are escalated when suspicious activity is noticed
			\item Controls are double-checked to anticipate any issues
			\item Individuals communicate with peers for knowledge sharing
			\item Employees are self-motivated to learn about costs and benefits
		\end{enumerate}
	\end{frame}
	\begin{frame}{D1-Q24: Answer}
		\begin{block}{Correct Answer: A}
			Management benefits most from an escalation process because risk and incidents are reported in a timely manner. Escalation posture among employees is best developed through training and awareness programs.
		\end{block}
	\end{frame}
	
	% ---- D1 Q25 ----
	\begin{frame}{D1-Q25: Main Objective of IT Risk Management}
		\begin{block}{Question}
			The MAIN objective of IT risk management is to:
		\end{block}
		\begin{enumerate}[(A)]
			\item Prevent loss of IT assets
			\item Provide timely management reports
			\item Ensure regulatory compliance
			\item Enable risk-aware business decisions
		\end{enumerate}
	\end{frame}
	\begin{frame}{D1-Q25: Answer}
		\begin{block}{Correct Answer: D}
			IT risk management should be conducted as part of enterprise-wide risk management, whose ultimate objective is to support risk-aware business decisions. Asset protection, timely reporting, and compliance are subordinate goals.
		\end{block}
	\end{frame}
	
	% ---- D1 Q26 ----
	\begin{frame}{D1-Q26: Multiple Governmental Jurisdictions}
		\begin{block}{Question}
			Which is the BEST approach for a global enterprise subject to regulation by multiple governmental jurisdictions with differing requirements?
		\end{block}
		\begin{enumerate}[(A)]
			\item Bring all locations into conformity with the aggregate requirements of all jurisdictions
			\item Bring all locations into conformity with generally accepted industry good practices
			\item Establish a baseline standard incorporating the requirements all jurisdictions have in common
			\item Establish baseline standards for all locations and add supplemental standards as required
		\end{enumerate}
	\end{frame}
	\begin{frame}{D1-Q26: Answer}
		\begin{block}{Correct Answer: D}
			It is more efficient to establish a baseline standard and then develop additional standards for locations that must meet specific requirements. Aggregate conformity places undue burden; good practices may not meet compliance; lowest common denominator risks noncompliance.
		\end{block}
	\end{frame}
	
	% ---- D1 Q27 ----
	\begin{frame}{D1-Q27: Final Sign-off on Risk Plan}
		\begin{block}{Question}
			Who MUST give final sign-off on the IT risk management plan?
		\end{block}
		\begin{enumerate}[(A)]
			\item IT auditors performing the risk assessment
			\item Business process owners
			\item Senior managers
			\item IT security administrators
		\end{enumerate}
	\end{frame}
	\begin{frame}{D1-Q27: Answer}
		\begin{block}{Correct Answer: C --- Senior managers}
			Senior managers understand performance metrics, approved the policies and standards that govern the enterprise, and have final responsibility for risk associated with audit findings and recommendations.
		\end{block}
	\end{frame}
	
	% ---- D1 Q28 ----
	\begin{frame}{D1-Q28: Security Boundary Before Assessment}
		\begin{block}{Question}
			Which is the PRIMARY reason a risk practitioner determines the security boundary prior to conducting a risk assessment?
		\end{block}
		\begin{enumerate}[(A)]
			\item To decide which laws and regulations apply
			\item To identify the scope of the risk assessment
			\item To identify the business owners of the system
			\item To decide whether quantitative or qualitative analysis is appropriate
		\end{enumerate}
	\end{frame}
	\begin{frame}{D1-Q28: Answer}
		\begin{block}{Correct Answer: B --- Identify the scope}
			Identifying the security boundary establishes the fundamental scope of inquiry, including which systems and components are subject to assessment. The boundary subsequently informs which laws apply and which business owners to consult.
		\end{block}
	\end{frame}
	
	% ---- D1 Q29 ----
	\begin{frame}{D1-Q29: Who Manages/Executes Risk Program}
		\begin{block}{Question}
			Which group would be MOST effective in managing and executing an enterprise's risk program?
		\end{block}
		\begin{enumerate}[(A)]
			\item Mid-level managers
			\item Senior managers
			\item Frontline employees
			\item The incident response team
		\end{enumerate}
	\end{frame}
	\begin{frame}{D1-Q29: Answer}
		\begin{block}{Correct Answer: A --- Mid-level managers}
			Mid-level managers are most centrally located within the organizational hierarchy, combining sufficient breadth of influence with proximity to day-to-day operations. Senior managers sponsor but are too high to manage; frontline employees lack influence.
		\end{block}
	\end{frame}
	
	% ---- D1 Q30 ----
	\begin{frame}{D1-Q30: IT Steering Committee Representation}
		\begin{block}{Question}
			The IT steering committee of a one-year-old start-up makes all IT decisions, including technology budget. It will be BEST represented by:
		\end{block}
		\begin{enumerate}[(A)]
			\item Members of the executive board
			\item High-level members of the IT department
			\item IT experts from outside the enterprise
			\item Key members from each department
		\end{enumerate}
	\end{frame}
	\begin{frame}{D1-Q30: Answer}
		\begin{block}{Correct Answer: D}
			The IT steering committee should be comprised of individuals from each department to ensure the entire enterprise is represented and all business objectives are more likely to be met.
		\end{block}
	\end{frame}
	
	% ---- D1 Q31 ----
	\begin{frame}{D1-Q31: First Line of Defense Purpose}
		\begin{block}{Question}
			Which is one of the MAIN purposes of the first line of defense in the three lines of defense model?
		\end{block}
		\begin{enumerate}[(A)]
			\item Ensure that financial controls are in place
			\item Ensure control deficiencies are addressed
			\item Ensure risk management practices are effective
			\item Ensure compliance with rules and regulations
		\end{enumerate}
	\end{frame}
	\begin{frame}{D1-Q31: Answer}
		\begin{block}{Correct Answer: B}
			In the first line, operational managers own the risk and address it by mitigating control deficiencies. Risk management is part of the second line; ensuring compliance is part of the second line.
		\end{block}
	\end{frame}
	
	% ---- D1 Q32 ----
	\begin{frame}{D1-Q32: Selecting Risk Response Technique}
		\begin{block}{Question}
			The PRIMARY consideration when selecting a risk response technique is:
		\end{block}
		\begin{enumerate}[(A)]
			\item Coverage of all identified risk
			\item Availability of resources
			\item Enterprise goals and objectives
			\item Standards and industry good practices
		\end{enumerate}
	\end{frame}
	\begin{frame}{D1-Q32: Answer}
		\begin{block}{Correct Answer: C}
			The risk response will be based primarily on the goals and objectives of the enterprise. Risk can harm these goals and must be mitigated according to priority. Standards and good practices are followed after the technique is selected.
		\end{block}
	\end{frame}
	
	% ---- D1 Q33 ----
	\begin{frame}{D1-Q33: First Policy for Start-up}
		\begin{block}{Question}
			The board of a one-year-old start-up asked the CIO to create all IT policies and procedures. Which should the CIO create FIRST?
		\end{block}
		\begin{enumerate}[(A)]
			\item The strategic IT plan
			\item The data classification scheme
			\item The information architecture document
			\item The technology infrastructure plan
		\end{enumerate}
	\end{frame}
	\begin{frame}{D1-Q33: Answer}
		\begin{block}{Correct Answer: A --- Strategic IT plan}
			The strategic IT plan is the first policy to create when developing an enterprise's governance model. Data classification, information architecture, and technology infrastructure plan are created afterward.
		\end{block}
	\end{frame}
	
	% ---- D1 Q34 ----
	\begin{frame}{D1-Q34: Most Concerning to Risk Practitioner}
		\begin{block}{Question}
			Which should be of MOST concern to a risk practitioner?
		\end{block}
		\begin{enumerate}[(A)]
			\item Failure to notify the public of an intrusion
			\item Failure to notify the police of an attempted intrusion
			\item Failure to internally report a successful attack
			\item Failure to examine access rights periodically
		\end{enumerate}
	\end{frame}
	\begin{frame}{D1-Q34: Answer}
		\begin{block}{Correct Answer: C}
			Failure to report a successful intrusion is a serious concern and could in some instances be interpreted as abetting. Public notification is not a requirement; attempted intrusions rarely require police notification.
		\end{block}
	\end{frame}
	
	% ---- D1 Q35 ----
	\begin{frame}{D1-Q35: Insight into Risk Capabilities}
		\begin{block}{Question}
			Which review will provide the MOST insight into an enterprise's risk management capabilities?
		\end{block}
		\begin{enumerate}[(A)]
			\item A capability maturity model review
			\item A capability comparison with industry standards or regulations
			\item A self-assessment of capabilities
			\item An internal audit review of capabilities
		\end{enumerate}
	\end{frame}
	\begin{frame}{D1-Q35: Answer}
		\begin{block}{Correct Answer: A}
			Capability maturity modeling allows an enterprise to understand its level of maturity in its risk capabilities, which is an indicator of operational readiness and effectiveness. The other options indicate existence or non-existence only.
		\end{block}
	\end{frame}
	
	% ---- D1 Q36 ----
	\begin{frame}{D1-Q36: Leveraging Internal Audit in ERM}
		\begin{block}{Question}
			As part of an ERM program, a risk practitioner BEST leverages the work performed by internal audit by having it:
		\end{block}
		\begin{enumerate}[(A)]
			\item Design, implement, and maintain the ERM process
			\item Manage and assess the overall risk awareness
			\item Evaluate ongoing changes to enterprise risk factors
			\item Assist in monitoring, evaluating, examining, and reporting on controls
		\end{enumerate}
	\end{frame}
	\begin{frame}{D1-Q36: Answer}
		\begin{block}{Correct Answer: D}
			The internal audit function is responsible for assisting management and the board in monitoring, evaluating, examining, and reporting on internal controls, regardless of whether an ERM function has been implemented.
		\end{block}
	\end{frame}
	
	% ---- D1 Q37 ----
	\begin{frame}{D1-Q37: Interconnections with External Systems}
		\begin{block}{Question}
			Which is the BIGGEST concern for a CISO regarding interconnections with systems outside the enterprise?
		\end{block}
		\begin{enumerate}[(A)]
			\item Requirements to comply with each other's contractual security obligations
			\item Uncertainty that the other system will be available as needed
			\item The ability to perform risk assessments on the other system
			\item Ensuring communications are encrypted through a VPN tunnel
		\end{enumerate}
	\end{frame}
	\begin{frame}{D1-Q37: Answer}
		\begin{block}{Correct Answer: A}
			Ensuring both systems comply with mutual contractual security obligations should be the primary concern. If one system fails to comply, both will likely miss their respective security obligations.
		\end{block}
	\end{frame}
	
	% ---- D1 Q38 ----
	\begin{frame}{D1-Q38: IT Organizational Structure}
		\begin{block}{Question}
			The CIO creates all IT policies and procedures, managed and approved by the IT steering committee, which makes all IT decisions including the technology budget. Which type of IT organizational structure is this?
		\end{block}
		\begin{enumerate}[(A)]
			\item Project-based
			\item Centralized
			\item Decentralized
			\item Divisional
		\end{enumerate}
	\end{frame}
	\begin{frame}{D1-Q38: Answer}
		\begin{block}{Correct Answer: B --- Centralized}
			Within a centralized IT organizational structure, one group makes all decisions for the entire enterprise.
		\end{block}
	\end{frame}
	
	% ---- D1 Q39 ----
	\begin{frame}{D1-Q39: Communicating Continuous Monitoring}
		\begin{block}{Question}
			The aggregated results of continuous monitoring activities are BEST communicated to:
		\end{block}
		\begin{enumerate}[(A)]
			\item The risk owner
			\item Technical staff
			\item The audit department
			\item The information security manager
		\end{enumerate}
	\end{frame}
	\begin{frame}{D1-Q39: Answer}
		\begin{block}{Correct Answer: A --- The risk owner}
			The risk owner is accountable for the fact that appropriate risk responses are executed in alignment with the enterprise's risk appetite. Technical staff benefit from granular data; audit and security managers are secondary audiences.
		\end{block}
	\end{frame}
	
	% ---- D1 Q40 ----
	\begin{frame}{D1-Q40: IT Risk Awareness Program}
		\begin{block}{Question}
			Which is a PRIMARY consideration when developing an IT risk awareness program?
		\end{block}
		\begin{enumerate}[(A)]
			\item Why technology risk is owned by IT
			\item How technology risk can affect each attendee's area of business
			\item How business process owners can transfer technology risk
			\item Why technology risk is more difficult to manage than other risk
		\end{enumerate}
	\end{frame}
	\begin{frame}{D1-Q40: Answer}
		\begin{block}{Correct Answer: B}
			Stakeholders must understand how IT-related risk affects the overall business and how potential threats and vulnerabilities can jeopardize the enterprise's people, processes, and technology. IT does not own technology risk.
		\end{block}
	\end{frame}
	
	% ---- D1 Q41 ----
	\begin{frame}{D1-Q41: IT Delivering Against Requirements}
		\begin{block}{Question}
			Management wants to ensure IT is successful in delivering against business requirements. Which BEST supports that effort?
		\end{block}
		\begin{enumerate}[(A)]
			\item An internal control system or framework
			\item A cost-benefit analysis
			\item A return on investment analysis
			\item A benchmark process
		\end{enumerate}
	\end{frame}
	\begin{frame}{D1-Q41: Answer}
		\begin{block}{Correct Answer: A}
			For IT to be successful in delivering against business requirements, management should develop an internal control system that supports its business requirements. Cost-benefit and ROI are metrics; benchmarking comes after a sound control framework.
		\end{block}
	\end{frame}
	
	% ---- D1 Q42 ----
	\begin{frame}{D1-Q42: Aligning Risk Appetite}
		\begin{block}{Question}
			It is MOST important that risk appetite be aligned with business objectives to ensure that:
		\end{block}
		\begin{enumerate}[(A)]
			\item Resources are directed toward areas of low risk tolerance
			\item Major risk is identified and eliminated
			\item IT and business goals are aligned
			\item The risk strategy is adequately communicated
		\end{enumerate}
	\end{frame}
	\begin{frame}{D1-Q42: Answer}
		\begin{block}{Correct Answer: A}
			Aligning risk appetite with business objectives allows an enterprise to evaluate and deploy valuable resources toward those objectives with low risk tolerance for loss. Risk cannot be eliminated; alignment is broader than IT/business; communication does not depend on alignment.
		\end{block}
	\end{frame}
	
	% ---- D1 Q43 ----
	\begin{frame}{D1-Q43: Out-of-Date Security Policy}
		\begin{block}{Question}
			Which causes the GREATEST concern to a risk practitioner reviewing a corporate information security policy that is out of date? The policy:
		\end{block}
		\begin{enumerate}[(A)]
			\item Was not reviewed within the last three years
			\item Is missing newer technologies/platforms
			\item Was not updated to account for new locations
			\item Does not enforce control monitoring
		\end{enumerate}
	\end{frame}
	\begin{frame}{D1-Q43: Answer}
		\begin{block}{Correct Answer: A}
			Not reviewing the policy for three years and updating it as necessary does not follow good practices and is the greatest concern. Policies are written at a level that does not require modification for specific technologies or locations.
		\end{block}
	\end{frame}
	
	% ---- D1 Q44 ----
	\begin{frame}{D1-Q44: Accountable for IT Business Risk}
		\begin{block}{Question}
			Who is accountable for business risk related to IT?
		\end{block}
		\begin{enumerate}[(A)]
			\item Chief information officer
			\item Chief financial officer
			\item Users of IT services
			\item Chief architect
		\end{enumerate}
	\end{frame}
	\begin{frame}{D1-Q44: Answer}
		\begin{block}{Correct Answer: C --- Users of IT services}
			The enterprise (i.e., the users of IT services) ultimately owns business-related risk, including the risk related to the use of IT. The business sets the mandate, provides resources, and monitors whether risk is being managed.
		\end{block}
	\end{frame}
	
	% ---- D1 Q45 ----
	\begin{frame}{D1-Q45: Expanding to Nearby Domestic Locations}
		\begin{block}{Question}
			An enterprise is expanding into new nearby domestic locations. Which is MOST important for a risk practitioner to report on?
		\end{block}
		\begin{enumerate}[(A)]
			\item Competitor analysis
			\item Legal and regulatory requirements
			\item Political issues
			\item Potential of natural disasters
		\end{enumerate}
	\end{frame}
	\begin{frame}{D1-Q45: Answer}
		\begin{block}{Correct Answer: B}
			Legal and regulatory requirements are most likely to change when moving to a nearby location because each municipality may enforce significantly different regulations, including environmental requirements, taxation, and others.
		\end{block}
	\end{frame}
	
	% ---- D1 Q46 ----
	\begin{frame}{D1-Q46: True Statement About IT Risk}
		\begin{block}{Question}
			Which of the following is true about IT risk?
		\end{block}
		\begin{enumerate}[(A)]
			\item IT risk cannot be assessed and measured quantitatively
			\item IT risk should be calculated separately from business risk
			\item IT risk management is the responsibility of the IT department
			\item IT risk exists regardless of whether it is detected or recognized
		\end{enumerate}
	\end{frame}
	\begin{frame}{D1-Q46: Answer}
		\begin{block}{Correct Answer: D}
			The enterprise must identify, acknowledge, and respond to risk; ignorance of risk is not acceptable. IT risk can be measured both quantitatively and qualitatively; it is one type of business risk; and it is the responsibility of senior management.
		\end{block}
	\end{frame}
	
	% ---- D1 Q47 ----
	\begin{frame}{D1-Q47: Selecting Risk Management Methodology}
		\begin{block}{Question}
			Which is MOST important when selecting an appropriate risk management methodology?
		\end{block}
		\begin{enumerate}[(A)]
			\item Risk culture
			\item Countermeasure analysis
			\item Cost-benefit analysis
			\item Risk transfer strategy
		\end{enumerate}
	\end{frame}
	\begin{frame}{D1-Q47: Answer}
		\begin{block}{Correct Answer: A --- Risk culture}
			Without understanding risk culture---how and why an enterprise makes decisions regarding risk---one cannot select a risk management methodology. Countermeasure and cost-benefit analysis target specific controls; not all risk can be transferred.
		\end{block}
	\end{frame}
	
	% ---- D1 Q48 ----
	\begin{frame}{D1-Q48: Compliance with Risk Appetite}
		\begin{block}{Question}
			Which BEST determines compliance with the risk appetite of an enterprise?
		\end{block}
		\begin{enumerate}[(A)]
			\item Balance between preventive and detective controls
			\item Inherent risk and acceptable risk level
			\item Residual risk and acceptable risk level
			\item Balance between countermeasures to threats and preventive controls
		\end{enumerate}
	\end{frame}
	\begin{frame}{D1-Q48: Answer}
		\begin{block}{Correct Answer: C}
			Considering residual risk in terms of acceptable risk yields risk that is appropriately balanced after the application of controls. Management can then decide to accept risk or apply additional controls.
		\end{block}
	\end{frame}
	
	% ---- D1 Q49 ----
	\begin{frame}{D1-Q49: Monitoring Control Effectiveness vs External Requirements}
		\begin{block}{Question}
			A key objective when monitoring information systems control effectiveness against the enterprise's external requirements is to:
		\end{block}
		\begin{enumerate}[(A)]
			\item Design the applicable information security controls for external audits
			\item Create the enterprise's information security policy provisions for third parties
			\item Ensure that the enterprise's legal obligations have been satisfied
			\item Identify those legal obligations that apply to the enterprise's security practices
		\end{enumerate}
	\end{frame}
	\begin{frame}{D1-Q49: Answer}
		\begin{block}{Correct Answer: C}
			Legal obligations are one of the principal external requirements that necessitate compliance monitoring. Control design occurs in risk treatment; policy creation occurs well before monitoring; obligations are identified before treatment.
		\end{block}
	\end{frame}
	
	% ---- D1 Q50 ----
	\begin{frame}{D1-Q50: Overall Risk Status}
		\begin{block}{Question}
			Risk monitoring provides timely information on the actual status of the enterprise with regard to risk. Which choice provides an overall risk status of the enterprise?
		\end{block}
		\begin{enumerate}[(A)]
			\item Risk management
			\item Risk analysis
			\item Risk appetite
			\item Risk profile
		\end{enumerate}
	\end{frame}
	\begin{frame}{D1-Q50: Answer}
		\begin{block}{Correct Answer: D --- Risk profile}
			The risk profile provides the current overall portfolio of identified risk to which the enterprise is exposed. Because the profile is updated with evolving and new risk, it provides the enterprise's current risk status.
		\end{block}
	\end{frame}
	
	% ---- D1 Q51 ----
	\begin{frame}{D1-Q51: Firewall Deviation from Security Policy}
		\begin{block}{Question}
			Which activity should a risk professional perform to determine whether firewall deployments are deviating from the enterprise's information security policy?
		\end{block}
		\begin{enumerate}[(A)]
			\item Review the firewall parameter settings
			\item Review the firewall intrusion prevention system logs
			\item Review the firewall hardening procedures
			\item Analyze the firewall log file for recent attacks
		\end{enumerate}
	\end{frame}
	\begin{frame}{D1-Q51: Answer}
		\begin{block}{Correct Answer: A}
			Firewall parameter settings tie into configurations linked to the governing security policy. If parameter settings differ from what policy states or requires, there is a deviation.
		\end{block}
	\end{frame}
	
	% ---- D1 Q52 ----
	\begin{frame}{D1-Q52: Continuous Risk Management Benefit}
		\begin{block}{Question}
			Which is the MOST significant benefit of implementing a continuous risk management process in an enterprise?
		\end{block}
		\begin{enumerate}[(A)]
			\item It complies with external regulators' requirements
			\item It captures changes to the enterprise's risk profile
			\item It ensures stakeholders are linked to the identified risk
			\item It ensures timely action is taken to mitigate risk
		\end{enumerate}
	\end{frame}
	\begin{frame}{D1-Q52: Answer}
		\begin{block}{Correct Answer: B}
			A continuous risk management process ensures changes to the enterprise's risk profile are captured. The other choices are secondary or inaccurate.
		\end{block}
	\end{frame}
	
	% ---- D1 Q53 ----
	\begin{frame}{D1-Q53: High Risk vs Moderate Appetite}
		\begin{block}{Question}
			Senior management defined the enterprise risk appetite as moderate. A business-critical application poses a high risk. What is the NEXT action?
		\end{block}
		\begin{enumerate}[(A)]
			\item Remove the high-risk application and replace it with another system
			\item Recommend that management increase its acceptable risk level
			\item Assess the impact of planned controls to the application
			\item Restrict access to the application only to trusted users
		\end{enumerate}
	\end{frame}
	\begin{frame}{D1-Q53: Answer}
		\begin{block}{Correct Answer: C}
			The risk practitioner should determine whether planned implementation of new controls may lower the risk from high to moderate or low before taking further action. Removing the app risks downtime; recommending a tolerance increase is the business owner's call.
		\end{block}
	\end{frame}
	
	% ---- D1 Q54 ----
	\begin{frame}{D1-Q54: Developing the Risk Profile}
		\begin{block}{Question}
			Which BEST assists in the development of the risk profile?
		\end{block}
		\begin{enumerate}[(A)]
			\item The presence of preventive and detective controls
			\item Inherent risk and detection risk
			\item Cost-benefit analysis of controls
			\item Likelihood and impact of risk
		\end{enumerate}
	\end{frame}
	\begin{frame}{D1-Q54: Answer}
		\begin{block}{Correct Answer: D}
			Likelihood and impact of risk help in the development of the risk profile. Controls alone, inherent/detection risk, and cost-benefit analysis do not develop the profile itself.
		\end{block}
	\end{frame}
	
	% ---- D1 Q55 ----
	\begin{frame}{D1-Q55: Employee Handbook in Multiple Regions}
		\begin{block}{Question}
			An enterprise expanded into Europe, Asia, and Latin America and has a single-version, multiple-language handbook last updated three years ago. Which is of MOST concern?
		\end{block}
		\begin{enumerate}[(A)]
			\item The handbook may not have been correctly translated
			\item Newer policies may not be included
			\item Expired policies may be included
			\item The handbook may violate local laws and regulations
		\end{enumerate}
	\end{frame}
	\begin{frame}{D1-Q55: Answer}
		\begin{block}{Correct Answer: D}
			Because customs and laws affect an enterprise's ability to operate in a location and vary by state and country, it is critical for the employee handbook to acknowledge and account for regional and national differences.
		\end{block}
	\end{frame}
	
	% ============================================================
	% DOMAIN 2 --- IT RISK ASSESSMENT (20%) --- Q56--Q100
	% ============================================================
	\section{Domain 2: IT Risk Assessment}
	
	% ---- D2 Q56 ----
	\begin{frame}{D2-Q56: Annual Loss Expectancy Formula}
		\begin{block}{Question}
			An asset's annual loss expectancy is calculated as the:
		\end{block}
		\begin{enumerate}[(A)]
			\item Exposure factor (EF) multiplied by the annualized rate of occurrence (ARO)
			\item Single loss expectancy (SLE) multiplied by the EF
			\item SLE multiplied by the ARO
			\item Asset value multiplied by the SLE
		\end{enumerate}
	\end{frame}
	\begin{frame}{D2-Q56: Answer}
		\begin{block}{Correct Answer: C --- ALE = SLE × ARO}
			ALE is calculated by multiplying the single loss expectancy (SLE) by the annualized rate of occurrence (ARO), the number of times the enterprise expects the loss to occur.
		\end{block}
	\end{frame}
	
	% ---- D2 Q57 ----
	\begin{frame}{D2-Q57: Effective Information Risk Program}
		\begin{block}{Question}
			Which is the BEST indicator of an effective information risk management program?
		\end{block}
		\begin{enumerate}[(A)]
			\item The security policy is made widely available
			\item Risk is considered before all decisions
			\item Security procedures are updated annually
			\item Risk assessments occur on an annual basis
		\end{enumerate}
	\end{frame}
	\begin{frame}{D2-Q57: Answer}
		\begin{block}{Correct Answer: B}
			Defining information risk in advance of business decisions best ensures that risk tolerance remains at approved levels. Policy availability, annual procedures updates, and annual assessments facilitate but are not as critical.
		\end{block}
	\end{frame}
	
	% ---- D2 Q58 ----
	\begin{frame}{D2-Q58: Risk Management Strategic Plans}
		\begin{block}{Question}
			Risk management strategic plans are MOST effective when developed for:
		\end{block}
		\begin{enumerate}[(A)]
			\item The enterprise as a whole
			\item Each individual system based on technology used
			\item Every location based on geographic threats
			\item End-to-end business processes
		\end{enumerate}
	\end{frame}
	\begin{frame}{D2-Q58: Answer}
		\begin{block}{Correct Answer: A}
			Risk management strategic plans are most effective when they are created and followed by the entire enterprise. Per-technology, per-location, or per-process plans create complexity and conflicting policies.
		\end{block}
	\end{frame}
	
	% ---- D2 Q59 ----
	\begin{frame}{D2-Q59: Considering Risk Appetite}
		\begin{block}{Question}
			Which is MOST important when considering the risk appetite of an enterprise?
		\end{block}
		\begin{enumerate}[(A)]
			\item The capacity of the enterprise to absorb loss
			\item The definition of responsibilities for risk management
			\item The line of business and the typical risk of the industry
			\item The culture and predisposition toward risk taking
		\end{enumerate}
	\end{frame}
	\begin{frame}{D2-Q59: Answer}
		\begin{block}{Correct Answer: D}
			When considering risk appetite, two major factors are relevant: the management culture and the predisposition toward risk taking. Capacity to absorb loss is important but does not influence appetite as much as culture.
		\end{block}
	\end{frame}
	
	% ---- D2 Q60 ----
	\begin{frame}{D2-Q60: Primary Purpose of ERM Framework}
		\begin{block}{Question}
			The PRIMARY purpose of adopting an enterprise-wide risk management framework is to:
		\end{block}
		\begin{enumerate}[(A)]
			\item Allow flexibility to adjust the risk response strategy
			\item Centralize responsibility for the maintenance of the risk response program
			\item Enable a consistent approach to risk response throughout the enterprise
			\item Avoid higher costs for risk reduction and audit strategies
		\end{enumerate}
	\end{frame}
	\begin{frame}{D2-Q60: Answer}
		\begin{block}{Correct Answer: C}
			Enabling consistent risk response is a key objective of the risk management framework. Risk management is the responsibility of all individuals; accountability lies with senior management and the board.
		\end{block}
	\end{frame}
	
	% ---- D2 Q61 ----
	\begin{frame}{D2-Q61: 15\% Project Deviation}
		\begin{block}{Question}
			Projects frequently exceed timelines or budget by nearly 10 percent; management advises a deviation of 15 percent is acceptable. This is an example of:
		\end{block}
		\begin{enumerate}[(A)]
			\item Risk avoidance
			\item Risk tolerance
			\item Risk acceptance
			\item Risk mitigation
		\end{enumerate}
	\end{frame}
	\begin{frame}{D2-Q61: Answer}
		\begin{block}{Correct Answer: B --- Risk tolerance}
			Risk tolerance is the permissible deviation from declared risk appetite levels. Avoidance steers clear of the activity; acceptance is a decision not to act; mitigation manages risk through countermeasures.
		\end{block}
	\end{frame}
	
	% ---- D2 Q62 ----
	\begin{frame}{D2-Q62: Computing Annual Loss Exposure}
		\begin{block}{Question}
			Which is MOST useful when computing annual loss exposure?
		\end{block}
		\begin{enumerate}[(A)]
			\item The cost of existing controls
			\item The number of vulnerabilities
			\item The net present value of the asset
			\item The business value of the asset
		\end{enumerate}
	\end{frame}
	\begin{frame}{D2-Q62: Answer}
		\begin{block}{Correct Answer: D}
			Annual loss exposure is a function of the value of the information asset and the impact if a given potential risk should materialize. Control cost, vulnerability count, and NPV do not directly determine loss exposure.
		\end{block}
	\end{frame}
	
	% ---- D2 Q63 ----
	\begin{frame}{D2-Q63: GRC Tool Major Risk}
		\begin{block}{Question}
			Which is a MAJOR risk associated with the use of governance, risk, and compliance (GRC) tools?
		\end{block}
		\begin{enumerate}[(A)]
			\item Misinterpretation of the dashboard's output
			\item Incomplete or inaccurate coverage in the enterprise for GRC
			\item Obsolescence of content
			\item Complex integration of diverse requirements
		\end{enumerate}
	\end{frame}
	\begin{frame}{D2-Q63: Answer}
		\begin{block}{Correct Answer: C --- Obsolescence of content}
			A GRC application must be updated regularly with current regulations and policies. Obsolete content renders a GRC outdated. The technology team should refresh the unified compliance framework (UCF) file quarterly and implement processes to address changes.
		\end{block}
	\end{frame}
	
	% ---- D2 Q64 ----
	\begin{frame}{D2-Q64: Need to Review Risk Practices}
		\begin{block}{Question}
			Which signifies the need to review an enterprise's risk practices?
		\end{block}
		\begin{enumerate}[(A)]
			\item Business owners regularly challenge risk assessment findings
			\item Manufacturing assigns its own internal risk management roles
			\item The finance department finds exceptions during its yearly risk review
			\item Sales department risk procedures were last reviewed 11 months ago
		\end{enumerate}
	\end{frame}
	\begin{frame}{D2-Q64: Answer}
		\begin{block}{Correct Answer: A}
			An enterprise's risk management practices must be clearly understood and supported by business stakeholders. Business owners challenging risk assessment findings either do not support or do not understand them clearly.
		\end{block}
	\end{frame}
	
	% ---- D2 Q65 ----
	\begin{frame}{D2-Q65: Correct Info Not Received in Time}
		\begin{block}{Question}
			The correct information was not received by the necessary recipients in a suitable time to allow proper action. This can be categorized as:
		\end{block}
		\begin{enumerate}[(A)]
			\item Integrity risk
			\item Availability risk
			\item Access risk
			\item Relevance risk
		\end{enumerate}
	\end{frame}
	\begin{frame}{D2-Q65: Answer}
		\begin{block}{Correct Answer: D --- Relevance risk}
			Relevance risk is a composite form of business risk, requiring both integrity and availability to be addressed. Transmitting information to the necessary recipients in a timely manner also creates tension with access (security) risk.
		\end{block}
	\end{frame}
	
	% ---- D2 Q66 ----
	\begin{frame}{D2-Q66: Justifying Security Program}
		\begin{block}{Question}
			Which risk assessment output is MOST suitable to help justify an enterprise information security program?
		\end{block}
		\begin{enumerate}[(A)]
			\item An inventory of risk that may impact the enterprise
			\item Documented threats to the enterprise
			\item Evaluation of the consequences
			\item A list of appropriate controls for addressing risk
		\end{enumerate}
	\end{frame}
	\begin{frame}{D2-Q66: Answer}
		\begin{block}{Correct Answer: D}
			A list of information security controls corresponding to risk scenarios is one of the primary deliverables of the risk assessment exercise. It demonstrates due consideration of risk and applicable controls to address the risk.
		\end{block}
	\end{frame}
	
	% ---- D2 Q67 ----
	\begin{frame}{D2-Q67: Ultimate Accountability for Risk}
		\begin{block}{Question}
			Accountability for risk ultimately belongs to the:
		\end{block}
		\begin{enumerate}[(A)]
			\item Chief risk officer
			\item Compliance officer
			\item Chief financial officer
			\item Board of directors
		\end{enumerate}
	\end{frame}
	\begin{frame}{D2-Q67: Answer}
		\begin{block}{Correct Answer: D --- Board of directors}
			The board of directors has ultimate accountability to shareholders, customers, employees, and the general public. The CRO, compliance officer, and CFO have responsibilities but are not ultimately accountable.
		\end{block}
	\end{frame}
	
	% ---- D2 Q68 ----
	\begin{frame}{D2-Q68: Outdated IT Standards}
		\begin{block}{Question}
			A risk assessment reveals that many corporate IT standards have not been updated. The BEST course of action is to:
		\end{block}
		\begin{enumerate}[(A)]
			\item Review the standards against current requirements and determine their adequacy
			\item Determine that the standards should be updated annually
			\item Report that IT standards are adequate and do not need updating
			\item Review the IT policy document to see how frequently standards should be updated
		\end{enumerate}
	\end{frame}
	\begin{frame}{D2-Q68: Answer}
		\begin{block}{Correct Answer: A}
			The risk practitioner should verify that the standards are still adequate; standards that are lacking should be updated. The other options do not determine adequacy or relevance.
		\end{block}
	\end{frame}
	
	% ---- D2 Q69 ----
	\begin{frame}{D2-Q69: Accountable for Third-Party App Risk}
		\begin{block}{Question}
			Marketing procures a third-party application found to pose data privacy regulatory risk in some regions. Who will be accountable for the risk if implemented globally?
		\end{block}
		\begin{enumerate}[(A)]
			\item The IT department
			\item The data privacy officer
			\item The chief risk officer
			\item The marketing department
		\end{enumerate}
	\end{frame}
	\begin{frame}{D2-Q69: Answer}
		\begin{block}{Correct Answer: D --- Marketing department}
			Marketing is the business owner of the application and therefore must be accountable. According to COBIT, accountability rests with those who own the required resources and have authority to approve the execution and accept the outcome.
		\end{block}
	\end{frame}
	
	% ---- D2 Q70 ----
	\begin{frame}{D2-Q70: Owns Risk Treatment Decisions}
		\begin{block}{Question}
			A risk practitioner is conducting a risk assessment of the local office's network infrastructure. Who owns the risk treatment decisions?
		\end{block}
		\begin{enumerate}[(A)]
			\item Control owner
			\item Chief information officer
			\item IT security manager
			\item Senior management
		\end{enumerate}
	\end{frame}
	\begin{frame}{D2-Q70: Answer}
		\begin{block}{Correct Answer: D --- Senior management}
			Senior management ultimately owns the risk and would make risk treatment decisions. The control owner, CIO alone, and IT security manager do not own the business risk.
		\end{block}
	\end{frame}
	
	% ---- D2 Q71 ----
	\begin{frame}{D2-Q71: Responsible for Third-Party App Risk}
		\begin{block}{Question}
			Marketing procures a third-party application that poses data privacy regulatory risk. Which role will be responsible for the risk it poses to the business if implemented globally?
		\end{block}
		\begin{enumerate}[(A)]
			\item The marketing department
			\item The IT department
			\item The data privacy officer
			\item The chief risk officer
		\end{enumerate}
	\end{frame}
	\begin{frame}{D2-Q71: Answer}
		\begin{block}{Correct Answer: B --- IT department}
			The IT department is responsible for the risk posed by this application. IT policy states that no tool or application can be implemented within production without a risk assessment and mitigation of all risk to an acceptable level.
		\end{block}
	\end{frame}
	
	% ---- D2 Q72 ----
	\begin{frame}{D2-Q72: Supporting Encryption Recommendation}
		\begin{block}{Question}
			Which provides the GREATEST support to a risk practitioner recommending encryption of corporate laptops and removable media as a risk mitigation measure?
		\end{block}
		\begin{enumerate}[(A)]
			\item Benchmarking with peers
			\item Evaluating public reports on encryption algorithms
			\item Developing a business case
			\item Scanning unencrypted systems for vulnerabilities
		\end{enumerate}
	\end{frame}
	\begin{frame}{D2-Q72: Answer}
		\begin{block}{Correct Answer: C --- Developing a business case}
			A business case documents how encryption addresses risk and illustrates that risk losses can be reduced by recommended controls. Peers have different environments; public reports are generic; scanning unencrypted systems without justification does not support the recommendation.
		\end{block}
	\end{frame}
	
	% ---- D2 Q73 ----
	\begin{frame}{D2-Q73: Additional Stakeholder for Privacy Risk}
		\begin{block}{Question}
			Marketing procures a third-party application that poses data privacy regulatory risk in certain regions. Which additional stakeholder should be included in reporting the risk?
		\end{block}
		\begin{enumerate}[(A)]
			\item Chief privacy officer
			\item Third-party vendor
			\item Business continuity manager
			\item Audit manager
		\end{enumerate}
	\end{frame}
	\begin{frame}{D2-Q73: Answer}
		\begin{block}{Correct Answer: A --- Chief privacy officer}
			The chief privacy officer supports defining controls that will be implemented to address and mitigate data transfer and data privacy requirements. If the application must be used, the CPO becomes a critical stakeholder who must be informed.
		\end{block}
	\end{frame}
	
	% ---- D2 Q74 ----
	\begin{frame}{D2-Q74: Cloud Email Retention Contract}
		\begin{block}{Question}
			A company learned its cloud email provider was never contracted to back up messages, despite a three-year retention policy. Which would have BEST safeguarded the company?
		\end{block}
		\begin{enumerate}[(A)]
			\item Providing the contractor with the record retention policy up front
			\item Validating the company policies to the provider's contract
			\item Providing the contractor with the email retention policy up front
			\item Backing up the data on the company's internal network nightly
		\end{enumerate}
	\end{frame}
	\begin{frame}{D2-Q74: Answer}
		\begin{block}{Correct Answer: B}
			The initial review of third-party services should confirm that vendors are contractually required to enforce all internal policies, including record retention. Providing policies does not legally bind the third party.
		\end{block}
	\end{frame}
	
	% ---- D2 Q75 ----
	\begin{frame}{D2-Q75: No Policy for Personal Devices}
		\begin{block}{Question}
			If an enterprise does not have a formal policy regarding personal devices in the workplace, which should the risk practitioner recommend?
		\end{block}
		\begin{enumerate}[(A)]
			\item Introduce remote wipe functionality
			\item Implement an exception process based on appropriate approvals
			\item Update incident response procedures
			\item Create and maintain an inventory of personal devices
		\end{enumerate}
	\end{frame}
	\begin{frame}{D2-Q75: Answer}
		\begin{block}{Correct Answer: B}
			If employees can attach personal devices without a formal policy, employees should be required to gain approval from management, documented and reviewed in a security exception form based on the exception process.
		\end{block}
	\end{frame}
	
	% ---- D2 Q76 ----
	\begin{frame}{D2-Q76: Applying Risk Management}
		\begin{block}{Question}
			To be effective, risk management should be applied to:
		\end{block}
		\begin{enumerate}[(A)]
			\item The elements identified by a risk assessment
			\item Any area that exceeds acceptable risk levels
			\item All enterprise activities
			\item Only those areas that have potential impact
		\end{enumerate}
	\end{frame}
	\begin{frame}{D2-Q76: Answer}
		\begin{block}{Correct Answer: C}
			Although not all enterprise activities entail unacceptable risk, the practice of risk management is ideally applied to all enterprise activities. Risk management must be holistic, not limited to areas that exceed acceptable levels.
		\end{block}
	\end{frame}
	
	% ---- D2 Q77 ----
	\begin{frame}{D2-Q77: Basis for Corporate Security Policy}
		\begin{block}{Question}
			Development of corporate information security policy should PRIMARILY be based on:
		\end{block}
		\begin{enumerate}[(A)]
			\item Vulnerabilities
			\item Threats
			\item Assets
			\item Impacts
		\end{enumerate}
	\end{frame}
	\begin{frame}{D2-Q77: Answer}
		\begin{block}{Correct Answer: C --- Assets}
			The corporate information security policy is based on management's commitment to protect the assets of the enterprise (and relevant information of its business partners) from threats, risk, and exposures that could occur.
		\end{block}
	\end{frame}
	
	% ---- D2 Q78 ----
	\begin{frame}{D2-Q78: Risk-Aware Business Decisions}
		\begin{block}{Question}
			Which MOST enables risk-aware business decisions?
		\end{block}
		\begin{enumerate}[(A)]
			\item Robust information security policies
			\item Exchange of accurate and timely information
			\item Skilled risk management personnel
			\item Skilled risk management personnel
		\end{enumerate}
	\end{frame}
	\begin{frame}{D2-Q78: Answer}
		\begin{block}{Correct Answer: B}
			Robust exchange of information enables management to optimize risk-related decisions. Accuracy and timeliness of information are critical success factors; the flow should be bidirectional.
		\end{block}
	\end{frame}
	
	% ---- D2 Q79 ----
	\begin{frame}{D2-Q79: Initial Stages of Risk Program}
		\begin{block}{Question}
			Which requirement MUST be met during the initial stages of developing a risk management program?
		\end{block}
		\begin{enumerate}[(A)]
			\item Management establishes ownership of identified risk
			\item Information security policies and standards are established
			\item A management committee exists to provide program oversight
			\item The context and purpose of the program are defined
		\end{enumerate}
	\end{frame}
	\begin{frame}{D2-Q79: Answer}
		\begin{block}{Correct Answer: D}
			Initial requirements to determine the enterprise's purpose for creating an information security risk management program include determining the desired outcomes and defining objectives. Ownership, policies, and oversight come later.
		\end{block}
	\end{frame}
	
	% ---- D2 Q80 ----
	\begin{frame}{D2-Q80: Effective BYOD Policy}
		\begin{block}{Question}
			An enterprise implemented a corporate BYOD policy to reduce the risk of data leakage. Which approach MOST enables the policy to be effective?
		\end{block}
		\begin{enumerate}[(A)]
			\item Obtaining signed acceptance from users
			\item Educating users on acceptable and unacceptable practices
			\item Requiring users to read the policy and future updates
			\item Clearly stating disciplinary action for noncompliance
		\end{enumerate}
	\end{frame}
	\begin{frame}{D2-Q80: Answer}
		\begin{block}{Correct Answer: B}
			A BYOD policy that is approved and communicated will not be effective without proper training of all users. Signed acceptance does not imply understanding; reading and disciplinary statements are secondary.
		\end{block}
	\end{frame}
	
	% ---- D2 Q81 ----
	\begin{frame}{D2-Q81: New Regulation First Action}
		\begin{block}{Question}
			A new regulation for safeguarding information processed by a specific transaction type comes to an IT manager's attention. The manager should FIRST:
		\end{block}
		\begin{enumerate}[(A)]
			\item Meet with stakeholders to decide how to comply
			\item Analyze the key risk in the compliance process
			\item Update the existing security/privacy policy
			\item Assess whether existing controls meet the regulation
		\end{enumerate}
	\end{frame}
	\begin{frame}{D2-Q81: Answer}
		\begin{block}{Correct Answer: D}
			The first step is to understand the impact and requirements of the new regulation, which includes assessing how the enterprise will comply and to what extent existing controls support compliance. The risk practitioner should then assess any existing gaps.
		\end{block}
	\end{frame}
	
	% ---- D2 Q82 ----
	\begin{frame}{D2-Q82: Document Identifying Compliance Risk}
		\begin{block}{Question}
			Which document BEST identifies an enterprise's compliance risk and corrective actions in progress to meet regulatory requirements?
		\end{block}
		\begin{enumerate}[(A)]
			\item An internal audit report
			\item A risk register
			\item An external audit report
			\item A risk assessment report
		\end{enumerate}
	\end{frame}
	\begin{frame}{D2-Q82: Answer}
		\begin{block}{Correct Answer: B --- Risk register}
			A risk register provides a report of all current identified risk within an enterprise (including compliance risk) and tracks the status of corrective actions or exceptions.
		\end{block}
	\end{frame}
	
	% ---- D2 Q83 ----
	\begin{frame}{D2-Q83: Impact on Security Governance Models}
		\begin{block}{Question}
			Which will have the MOST significant impact on standard information security governance models?
		\end{block}
		\begin{enumerate}[(A)]
			\item Number of employees and consulting staff
			\item Cultural differences between physical locations
			\item Complexity of the organizational structure
			\item Uncertainty in legislative requirements
		\end{enumerate}
	\end{frame}
	\begin{frame}{D2-Q83: Answer}
		\begin{block}{Correct Answer: C}
			Information security governance models are highly dependent on the complexity of the organizational structure. Elements include multiple business units, dispersion of functions, multiple leadership hierarchies, and multiple lines of communication.
		\end{block}
	\end{frame}
	
	% ---- D2 Q84 ----
	\begin{frame}{D2-Q84: Highest Tolerance for Public Cloud}
		\begin{block}{Question}
			Senior management will MOST likely have the highest tolerance for moving which to a public cloud?
		\end{block}
		\begin{enumerate}[(A)]
			\item Credit-card processing
			\item Research and development
			\item The legacy financial system
			\item The corporate email system
		\end{enumerate}
	\end{frame}
	\begin{frame}{D2-Q84: Answer}
		\begin{block}{Correct Answer: D}
			Of the options, the corporate email system has the least competitive distinction, complexity, and sensitive/highly classified information. Moving processes to the cloud should consider criticality, complexity, and classification of the data.
		\end{block}
	\end{frame}
	
	% ---- D2 Q85 ----
	\begin{frame}{D2-Q85: Risk Profile Consideration}
		\begin{block}{Question}
			Which item is MOST important to consider in relation to a risk profile?
		\end{block}
		\begin{enumerate}[(A)]
			\item A summary of regional loss events
			\item Aggregated risk to the enterprise
			\item A description of critical risk
			\item An analysis of historical loss events
		\end{enumerate}
	\end{frame}
	\begin{frame}{D2-Q85: Answer}
		\begin{block}{Correct Answer: B}
			The risk profile is based on the aggregated risk to the enterprise, including historical risk, critical risk, and emerging risk. Regional loss events, critical risk alone, and historical loss analysis are subsets.
		\end{block}
	\end{frame}
	
	% ---- D2 Q86 ----
	\begin{frame}{D2-Q86: Collaborative Research Cloud}
		\begin{block}{Question}
			Which cloud deployment model is MOST appropriate for a collaborative research program between universities?
		\end{block}
		\begin{enumerate}[(A)]
			\item Private cloud
			\item Community cloud
			\item Public cloud
			\item Hybrid cloud
		\end{enumerate}
	\end{frame}
	\begin{frame}{D2-Q86: Answer}
		\begin{block}{Correct Answer: B --- Community cloud}
			A community cloud is shared between several enterprises and supports a specific community with a common mission or interest. Private is for a single enterprise; public is for the general public; hybrid combines models.
		\end{block}
	\end{frame}
	
	% ---- D2 Q87 ----
	\begin{frame}{D2-Q87: Control Environment Review Line}
		\begin{block}{Question}
			Which line of defense is MOST likely involved when conducting a review of the adequacy of the control environment?
		\end{block}
		\begin{enumerate}[(A)]
			\item The first line of defense
			\item The second line of defense
			\item The third line of defense
			\item Third-party line of defense
		\end{enumerate}
	\end{frame}
	\begin{frame}{D2-Q87: Answer}
		\begin{block}{Correct Answer: A --- First line}
			A control environment review is a first line of defense activity, where control owners are responsible for adequacy of the control environment. The second line provides risk-based consultation and monitoring; the third provides independent assurance.
		\end{block}
	\end{frame}
	
	% ---- D2 Q88 ----
	\begin{frame}{D2-Q88: Board Concern on Risk Capability}
		\begin{block}{Question}
			Which issue would be of MOST concern to the board when assessing risk management capability?
		\end{block}
		\begin{enumerate}[(A)]
			\item The third line acts independently of the other lines
			\item The second line is checking and challenging the first line
			\item Internal audit is preparing the risk management strategy
			\item Each line of defense is conducting its own planning independently
		\end{enumerate}
	\end{frame}
	\begin{frame}{D2-Q88: Answer}
		\begin{block}{Correct Answer: D}
			Joint planning across the three lines of defense is key to achieving an effective risk management capability. If each line is planning independently, the lines of defense are not working as expected.
		\end{block}
	\end{frame}
	
	% ---- D2 Q89 ----
	\begin{frame}{D2-Q89: Supporting Code of Ethics}
		\begin{block}{Question}
			Which MOST effectively supports adherence to an enterprise's code of ethics?
		\end{block}
		\begin{enumerate}[(A)]
			\item Ensuring periodic training, evaluation, and attestation of employees
			\item Performing background checks at hire
			\item Providing clear enterprise policies and standards
			\item Providing continuous awareness of ethical requirements
		\end{enumerate}
	\end{frame}
	\begin{frame}{D2-Q89: Answer}
		\begin{block}{Correct Answer: A}
			Attestation to comply with the code of conduct during ethics training ensures employees have a clear idea of what is expected. Training also provides evaluation to measure understanding. Background checks and policies alone are less effective.
		\end{block}
	\end{frame}
	
	% ---- D2 Q90 ----
	\begin{frame}{D2-Q90: Explaining Zero-Day Exploit}
		\begin{block}{Question}
			Who is responsible for explaining the ramifications of a new zero-day exploit to senior management?
		\end{block}
		\begin{enumerate}[(A)]
			\item Chief operating officer
			\item Chief risk officer
			\item Chief information security officer
			\item Chief information officer
		\end{enumerate}
	\end{frame}
	\begin{frame}{D2-Q90: Answer}
		\begin{block}{Correct Answer: B --- Chief risk officer}
			The chief risk officer is the most senior official accountable for all aspects of risk management across the enterprise, including explaining risk to senior management. The CISO leads the security program but is not responsible for enterprise risk acceptance.
		\end{block}
	\end{frame}
	
	% ---- D2 Q91 ----
	\begin{frame}{D2-Q91: Information Classification Responsibility}
		\begin{block}{Question}
			The application of information classification is the responsibility of the:
		\end{block}
		\begin{enumerate}[(A)]
			\item Information security officer
			\item Information owner
			\item Information systems auditor
			\item Information custodian
		\end{enumerate}
	\end{frame}
	\begin{frame}{D2-Q91: Answer}
		\begin{block}{Correct Answer: B --- Information owner}
			The information owner determines classification based on the criticality and sensitivity of information. The security officer, auditor, and custodian have different responsibilities.
		\end{block}
	\end{frame}
	
	% ---- D2 Q92 ----
	\begin{frame}{D2-Q92: Three Lines of Defense Benefit}
		\begin{block}{Question}
			The MOST significant benefit of using the three lines of defense model in a risk management framework is that it:
		\end{block}
		\begin{enumerate}[(A)]
			\item Ensures ongoing success of the risk management initiative
			\item Enhances communication between various stakeholders
			\item Clarifies essential roles of the key stakeholders
			\item Helps risk owners in selecting appropriate control owners
		\end{enumerate}
	\end{frame}
	\begin{frame}{D2-Q92: Answer}
		\begin{block}{Correct Answer: C}
			The three lines of defense model is designed to assign roles related to risk, controls, and security while better defining and understanding the relationships between them.
		\end{block}
	\end{frame}
	
	% ---- D2 Q93 ----
	\begin{frame}{D2-Q93: First Line of Defense Role}
		\begin{block}{Question}
			Which risk management role is part of the first line of defense?
		\end{block}
		\begin{enumerate}[(A)]
			\item Chief risk officer
			\item Risk steering committee
			\item Risk owner
			\item Board of directors
		\end{enumerate}
	\end{frame}
	\begin{frame}{D2-Q93: Answer}
		\begin{block}{Correct Answer: C --- Risk owner}
			The first line of defense is operational management, and risk owners are part of operational management. The CRO and risk steering committee are second line; the board provides direction and oversight.
		\end{block}
	\end{frame}
	
	% ---- D2 Q94 ----
	\begin{frame}{D2-Q94: Ultimate Objective of Risk Governance}
		\begin{block}{Question}
			What is the ULTIMATE objective of risk governance?
		\end{block}
		\begin{enumerate}[(A)]
			\item Benefits realization
			\item Risk optimization
			\item Resource optimization
			\item Value creation
		\end{enumerate}
	\end{frame}
	\begin{frame}{D2-Q94: Answer}
		\begin{block}{Correct Answer: D --- Value creation}
			Value creation is the main objective of risk governance and is achieved when the three underlying objectives (benefits realization, risk optimization, and resource optimization) are balanced.
		\end{block}
	\end{frame}
	
	% ---- D2 Q95 ----
	\begin{frame}{D2-Q95: Second Line of Defense Activity}
		\begin{block}{Question}
			Which activity would be a direct responsibility of the second line of defense?
		\end{block}
		\begin{enumerate}[(A)]
			\item Establish a risk management framework
			\item Identify, assess, and respond to risk
			\item Implement information security controls
			\item Test and report control effectiveness
		\end{enumerate}
	\end{frame}
	\begin{frame}{D2-Q95: Answer}
		\begin{block}{Correct Answer: A}
			Establishing a risk management framework, providing awareness training, and supervising overall risk management are responsibilities of the second line of defense. Identifying/assessing/responding is first line; testing controls is third line.
		\end{block}
	\end{frame}
	
	% ---- D2 Q96 ----
	\begin{frame}{D2-Q96: Ethical Behavior Contribution}
		\begin{block}{Question}
			Developing and practicing ethical behavior within an enterprise contributes the most to:
		\end{block}
		\begin{enumerate}[(A)]
			\item Determining the risk appetite
			\item Building the risk culture
			\item Enhancing risk management
			\item Increasing risk tolerance
		\end{enumerate}
	\end{frame}
	\begin{frame}{D2-Q96: Answer}
		\begin{block}{Correct Answer: B --- Building the risk culture}
			Practicing professional ethics within an enterprise helps stakeholders review risk associated with unethical behavior before making decisions. This establishes the risk culture of the enterprise.
		\end{block}
	\end{frame}
	
	% ---- D2 Q97 ----
	\begin{frame}{D2-Q97: Benefits of Professional Ethics}
		\begin{block}{Question}
			Which is the MOST important benefit when stakeholders promote and follow professional ethics?
		\end{block}
		\begin{enumerate}[(A)]
			\item It increases reliability in risk management processes
			\item It increases risk management effectiveness within an enterprise
			\item It improves compliance with regulatory requirements
			\item It improves the enterprise's reputation and business growth
		\end{enumerate}
	\end{frame}
	\begin{frame}{D2-Q97: Answer}
		\begin{block}{Correct Answer: B}
			When stakeholders follow ethical practices, actions and decisions are based on reviewing the results and impact of those actions. This promotes a risk-based approach to decision-making and enhances overall risk management.
		\end{block}
	\end{frame}
	
	% ---- D2 Q98 ----
	\begin{frame}{D2-Q98: Best Proactive Ethics Approach}
		\begin{block}{Question}
			Which is the BEST proactive approach for practicing professional ethics within an enterprise?
		\end{block}
		\begin{enumerate}[(A)]
			\item Develop an ethics policy
			\item Implement processes to report unethical behavior
			\item Define procedures for disciplinary action
			\item Provide ethics awareness training
		\end{enumerate}
	\end{frame}
	\begin{frame}{D2-Q98: Answer}
		\begin{block}{Correct Answer: D}
			Awareness training is the most proactive approach to improve professional ethics within an enterprise. Training informs employees on the code of ethics and enables them to make decisions and take appropriate action. Reporting and discipline are reactive.
		\end{block}
	\end{frame}
	
	% ---- D2 Q99 ----
	\begin{frame}{D2-Q99: Ensuring 3LoD Effectiveness}
		\begin{block}{Question}
			Which is the MOST critical to ensure the three lines of defense work together effectively?
		\end{block}
		\begin{enumerate}[(A)]
			\item A documented disaster recovery plan
			\item Automated operational processes
			\item Independent external audits
			\item Clear, detailed roles and responsibilities
		\end{enumerate}
	\end{frame}
	\begin{frame}{D2-Q99: Answer}
		\begin{block}{Correct Answer: D}
			Without clear, detailed job descriptions for employees working in each line of defense, there can be confusion on what is expected at each level, which may hinder effectiveness.
		\end{block}
	\end{frame}
	
	% ---- D2 Q100 ----
	\begin{frame}{D2-Q100: What Should Drive IT Plan}
		\begin{block}{Question}
			Which should drive the IT plan?
		\end{block}
		\begin{enumerate}[(A)]
			\item Strategic planning and business requirements
			\item Technology and operational procedures
			\item Compliance with laws and regulations
			\item Project plans and stakeholder requirements
		\end{enumerate}
	\end{frame}
	\begin{frame}{D2-Q100: Answer}
		\begin{block}{Correct Answer: A}
			IT exists to support business objectives. Management of enterprise IT should align the IT plan closely with the business. Compliance and project plans are evaluated in the context of business risk.
		\end{block}
	\end{frame}
	
	% ============================================================
	% DOMAIN 3 --- RISK RESPONSE AND REPORTING (32%) --- Q101--Q165
	% ============================================================
	\section{Domain 3: Risk Response and Reporting}
	
	% ---- D3 Q101 ----
	\begin{frame}{D3-Q101: Outsourced Noncore Process Concern}
		\begin{block}{Question}
			Which outcome of outsourcing noncore processes is of GREATEST concern to management?
		\end{block}
		\begin{enumerate}[(A)]
			\item Total cost of ownership exceeds projections
			\item Internal information systems experience is lost
			\item Employees of the vendor are disloyal to the client
			\item Processing of sensitive data is subcontracted by the vendor
		\end{enumerate}
	\end{frame}
	\begin{frame}{D3-Q101: Answer}
		\begin{block}{Correct Answer: D}
			The greatest risk in third-party relationships is ceding direct control of IS processes. Subcontracting increases this risk; the subcontracting process must be reviewed because sensitive data are involved.
		\end{block}
	\end{frame}
	
	% ---- D3 Q102 ----
	\begin{frame}{D3-Q102: Proposing Risk Treatment}
		\begin{block}{Question}
			When proposing the implementation of a specific risk treatment, a risk practitioner PRIMARILY uses a:
		\end{block}
		\begin{enumerate}[(A)]
			\item Technical evaluation report
			\item Business case
			\item Vulnerability assessment report
			\item Budgetary requirements
		\end{enumerate}
	\end{frame}
	\begin{frame}{D3-Q102: Answer}
		\begin{block}{Correct Answer: B --- Business case}
			A manager should base any proposed risk response on a risk assessment in the context of business objectives. The manager must document costs of controls and compare them against benefits of the risk response.
		\end{block}
	\end{frame}
	
	% ---- D3 Q103 ----
	\begin{frame}{D3-Q103: Most Important for Effective Risk Management}
		\begin{block}{Question}
			Which is MOST important for effective risk management?
		\end{block}
		\begin{enumerate}[(A)]
			\item Assignment of risk owners to identified risk
			\item Ensuring compliance with regulatory requirements
			\item Integration of risk management into operational processes
			\item Implementation of a risk avoidance strategy
		\end{enumerate}
	\end{frame}
	\begin{frame}{D3-Q103: Answer}
		\begin{block}{Correct Answer: A}
			It is of utmost importance to assign risk to individual owners and therein maximize accountability. Regulatory compliance is a relatively small part of risk management; risk avoidance is not always feasible.
		\end{block}
	\end{frame}
	
	% ---- D3 Q104 ----
	\begin{frame}{D3-Q104: Most Important Outsourcing Question}
		\begin{block}{Question}
			An enterprise outsourced several business functions including IT development, data hosting, and support. What is the MOST important question regarding the outsourcing arrangements?
		\end{block}
		\begin{enumerate}[(A)]
			\item Are policies and procedures in place to handle security exceptions?
			\item Is the supplier meeting the terms of the SLAs?
			\item Is the provider's security program based on an international standard?
			\item Are specific security controls mandated in the outsourcing contract?
		\end{enumerate}
	\end{frame}
	\begin{frame}{D3-Q104: Answer}
		\begin{block}{Correct Answer: D}
			Without enumerating security requirements directly in the outsourcing contract, the outsourcing company has no assurance that the provider will comply with specific security requirements.
		\end{block}
	\end{frame}
	
	% ---- D3 Q105 ----
	\begin{frame}{D3-Q105: Basis for Establishing Risk Ownership}
		\begin{block}{Question}
			Which activity provides the BEST basis for establishing risk ownership?
		\end{block}
		\begin{enumerate}[(A)]
			\item Documenting interdependencies between departments
			\item Mapping identified risk to a specific business process
			\item Referring to available RACI charts for business processes
			\item Distributing risk among identified business process owners
		\end{enumerate}
	\end{frame}
	\begin{frame}{D3-Q105: Answer}
		\begin{block}{Correct Answer: B}
			Mapping identified risk to a specific business process helps identify the process owner. Aggregation of related business processes results in identification of the prospective risk owner. Ownership cannot be shared.
		\end{block}
	\end{frame}
	
	% ---- D3 Q106 ----
	\begin{frame}{D3-Q106: Risk Response Report Recommendations}
		\begin{block}{Question}
			A risk response report includes recommendations for:
		\end{block}
		\begin{enumerate}[(A)]
			\item Acceptance
			\item Assessment
			\item Evaluation
			\item Quantification
		\end{enumerate}
	\end{frame}
	\begin{frame}{D3-Q106: Answer}
		\begin{block}{Correct Answer: A --- Acceptance}
			Acceptance of a risk is an alternative to be considered in the risk response process. Assessment, evaluation, and quantification occur earlier in the risk process.
		\end{block}
	\end{frame}
	
	% ---- D3 Q107 ----
	\begin{frame}{D3-Q107: Reduced to Achieve Acceptable Risk}
		\begin{block}{Question}
			Which is MOST likely to be reduced to achieve acceptable risk?
		\end{block}
		\begin{enumerate}[(A)]
			\item Risk appetite
			\item Control risk
			\item Residual risk
			\item Inherent risk
		\end{enumerate}
	\end{frame}
	\begin{frame}{D3-Q107: Answer}
		\begin{block}{Correct Answer: C --- Residual risk}
			Residual risk is the remaining risk after management has implemented a risk response. Acceptable risk is achieved when residual risk is reduced to levels within the enterprise's risk appetite.
		\end{block}
	\end{frame}
	
	% ---- D3 Q108 ----
	\begin{frame}{D3-Q108: Denial-of-Service Vulnerability Decision}
		\begin{block}{Question}
			A global financial institution decided not to take further action on a denial-of-service vulnerability. The MOST likely reason is:
		\end{block}
		\begin{enumerate}[(A)]
			\item The needed countermeasure is too complicated to deploy
			\item There are sufficient safeguards to prevent this risk
			\item The likelihood of the risk occurring is unknown
			\item The cost of countermeasure outweighs the value of the asset and potential loss
		\end{enumerate}
	\end{frame}
	\begin{frame}{D3-Q108: Answer}
		\begin{block}{Correct Answer: D}
			An enterprise may decide to accept a specific risk because the protection would cost more than the potential loss.
		\end{block}
	\end{frame}
	
	% ---- D3 Q109 ----
	\begin{frame}{D3-Q109: Two-Factor Authentication Cost-Benefit}
		\begin{block}{Question}
			Which is MOST relevant to include in a cost-benefit analysis of a two-factor authentication system?
		\end{block}
		\begin{enumerate}[(A)]
			\item The approved budget of the project
			\item The frequency of incidents
			\item The annual loss expectancy of security incidents
			\item The total cost of ownership
		\end{enumerate}
	\end{frame}
	\begin{frame}{D3-Q109: Answer}
		\begin{block}{Correct Answer: D}
			Total cost of ownership is the most relevant piece of information because it establishes a cost baseline that must be considered for the full life cycle of the control. Frequency and ALE have indirect relationships.
		\end{block}
	\end{frame}
	
	% ---- D3 Q110 ----
	\begin{frame}{D3-Q110: Most Important Outsourcing Contract Part}
		\begin{block}{Question}
			Which is the MOST important part of any outsourcing contract?
		\end{block}
		\begin{enumerate}[(A)]
			\item The right to audit the outsourcing provider
			\item Provisions to assess the compliance of the provider
			\item Procedures for dealing with incident notification
			\item Requirements to encrypt hosted data
		\end{enumerate}
	\end{frame}
	\begin{frame}{D3-Q110: Answer}
		\begin{block}{Correct Answer: B}
			If a contract contains no provision to monitor and hold a supplier accountable for security, then the outsourcing enterprise cannot ensure compliance or proper handling of its data.
		\end{block}
	\end{frame}
	
	% ---- D3 Q111 ----
	\begin{frame}{D3-Q111: Greatest Risk with Cloud Provider}
		\begin{block}{Question}
			Which poses the GREATEST risk to an enterprise that recently engaged a cloud provider?
		\end{block}
		\begin{enumerate}[(A)]
			\item The cloud provider's primary facility is in the same vicinity
			\item The service level agreement is ambiguous
			\item References from other customers were not obtained
			\item Auditing the vendor requires dependence on a third-party audit firm
		\end{enumerate}
	\end{frame}
	\begin{frame}{D3-Q111: Answer}
		\begin{block}{Correct Answer: B}
			If the service level agreement is ambiguous, it will be difficult to determine whether the provider complies. References cannot provide reasonable assurance; third-party audits are permitted under SSAE 18.
		\end{block}
	\end{frame}
	
	% ---- D3 Q112 ----
	\begin{frame}{D3-Q112: Insuring Against Unmitigable Risk}
		\begin{block}{Question}
			When a risk cannot be sufficiently mitigated through manual or automatic controls, which option will BEST protect the enterprise from the potential financial impact?
		\end{block}
		\begin{enumerate}[(A)]
			\item Insuring against the risk
			\item Updating the IT risk registry
			\item Improving staff training in the risk area
			\item Outsourcing the related process to a third party
		\end{enumerate}
	\end{frame}
	\begin{frame}{D3-Q112: Answer}
		\begin{block}{Correct Answer: A}
			An insurance policy can compensate the enterprise monetarily for the impact of the risk by transferring the risk to the insurance company. Updating the register changes perception only; outsourcing does not remove the risk.
		\end{block}
	\end{frame}
	
	% ---- D3 Q113 ----
	\begin{frame}{D3-Q113: Cost of Mitigation Exceeds Benefit}
		\begin{block}{Question}
			After a risk assessment, it is determined that the cost to mitigate the risk is much greater than the benefit to be derived. The BEST risk response is that the risk be:
		\end{block}
		\begin{enumerate}[(A)]
			\item Treated
			\item Terminated
			\item Accepted
			\item Transferred
		\end{enumerate}
	\end{frame}
	\begin{frame}{D3-Q113: Answer}
		\begin{block}{Correct Answer: C --- Accepted}
			When the cost of control is more than the cost of the potential impact, the risk should be accepted.
		\end{block}
	\end{frame}
	
	% ---- D3 Q114 ----
	\begin{frame}{D3-Q114: Improving Continuity of Operations}
		\begin{block}{Question}
			A data center provider's clients complain of frequent service disruptions. Management sets up a project team to improve continuity. Which risk management activity will MOST benefit this initiative?
		\end{block}
		\begin{enumerate}[(A)]
			\item Develop key risk indicators
			\item Determine risk and control ownership
			\item Update the risk register
			\item Develop IT risk scenarios
		\end{enumerate}
	\end{frame}
	\begin{frame}{D3-Q114: Answer}
		\begin{block}{Correct Answer: D}
			The relationship between IT risk scenarios and business impact must be established to understand the effects of possible adverse events on enterprise objectives. A risk scenario directly identifies circumstances that can adversely affect business or assets.
		\end{block}
	\end{frame}
	
	% ---- D3 Q115 ----
	\begin{frame}{D3-Q115: Reciprocal Agreement for Equipment}
		\begin{block}{Question}
			Critical IT equipment will be delivered late due to flooding, but a reciprocal agreement exists with another company for a replacement until it arrives. This is an example of risk:
		\end{block}
		\begin{enumerate}[(A)]
			\item Transfer
			\item Avoidance
			\item Acceptance
			\item Mitigation
		\end{enumerate}
	\end{frame}
	\begin{frame}{D3-Q115: Answer}
		\begin{block}{Correct Answer: D --- Mitigation}
			Risk mitigation attempts to reduce the impact when a risk event occurs. Making plans such as a reciprocal arrangement reduces the consequence of the risk event.
		\end{block}
	\end{frame}
	
	% ---- D3 Q116 ----
	\begin{frame}{D3-Q116: Selecting an Appropriate Risk Response}
		\begin{block}{Question}
			Which would BEST help an enterprise select an appropriate risk response?
		\end{block}
		\begin{enumerate}[(A)]
			\item An analysis of change in the risk environment
			\item An analysis of risk that can be transferred
			\item An analysis of likelihood and impact of risk scenarios
			\item An analysis of control costs and benefits
		\end{enumerate}
	\end{frame}
	\begin{frame}{D3-Q116: Answer}
		\begin{block}{Correct Answer: D}
			An analysis of costs and benefits for controls helps an enterprise understand if it can mitigate the risk to an acceptable level. Likelihood and impact establish the risk level but are inadequate for selecting the response.
		\end{block}
	\end{frame}
	
	% ---- D3 Q117 ----
	\begin{frame}{D3-Q117: Optimal Return on Security Investment}
		\begin{block}{Question}
			Which leads to the BEST optimal return on security investment?
		\end{block}
		\begin{enumerate}[(A)]
			\item Deploying maximum protection across all information assets
			\item Focusing on the most important information assets and then determining their protection
			\item Deploying minimum protection across all information assets
			\item Investing only after a major security incident is reported
		\end{enumerate}
	\end{frame}
	\begin{frame}{D3-Q117: Answer}
		\begin{block}{Correct Answer: B}
			To optimize return on security investment, focus on identifying important information assets and protect them appropriately. Maximum protection overprotects less critical assets; minimum protection compromises critical assets; reactive investing may be too late.
		\end{block}
	\end{frame}
	
	% ---- D3 Q118 ----
	\begin{frame}{D3-Q118: Compromised Administrator Account Impact}
		\begin{block}{Question}
			Which control can reduce the potential impact of a malicious hacker who gains access to an administrator account?
		\end{block}
		\begin{enumerate}[(A)]
			\item Multifactor authentication
			\item Audit logging
			\item Least privilege
			\item Password policy
		\end{enumerate}
	\end{frame}
	\begin{frame}{D3-Q118: Answer}
		\begin{block}{Correct Answer: C --- Least privilege}
			Least privilege can reduce the impact of a compromised account within the scope of its intended use and limit impact to the enterprise as a whole. MFA is already bypassed; audit logging is a lagging indicator; password policy does not limit scope.
		\end{block}
	\end{frame}
	
	% ---- D3 Q119 ----
	\begin{frame}{D3-Q119: Critical Element of Effective Risk Program}
		\begin{block}{Question}
			Which is critical to the risk practitioner for an effective risk management program?
		\end{block}
		\begin{enumerate}[(A)]
			\item Risk response strategy
			\item Risk register content
			\item Risk profile content
			\item Risk owners and accountability
		\end{enumerate}
	\end{frame}
	\begin{frame}{D3-Q119: Answer}
		\begin{block}{Correct Answer: D}
			The identification of risk owners is critical because risk owners must make informed and cost-effective decisions regarding appropriate controls to mitigate their owned risk. Strategy, register, and profile are tools that require proactive users with accountability.
		\end{block}
	\end{frame}
	
	% ---- D3 Q120 ----
	\begin{frame}{D3-Q120: Anti-Malware Cost Exceeds Loss}
		\begin{block}{Question}
			Where the cost of anti-malware exceeds the loss expectancy of malware threats, what is the MOST viable risk response?
		\end{block}
		\begin{enumerate}[(A)]
			\item Risk elimination
			\item Risk acceptance
			\item Risk transfer
			\item Risk mitigation
		\end{enumerate}
	\end{frame}
	\begin{frame}{D3-Q120: Answer}
		\begin{block}{Correct Answer: B --- Risk acceptance}
			When the cost of a risk response exceeds the loss expectancy, the most viable risk response is risk acceptance. Transfer and mitigation must also be based on cost-benefit analysis.
		\end{block}
	\end{frame}
	
	% ---- D3 Q121 ----
	\begin{frame}{D3-Q121: Best Example of Risk Avoidance}
		\begin{block}{Question}
			Which is the BEST example of risk avoidance behavior?
		\end{block}
		\begin{enumerate}[(A)]
			\item Taking no action against the risk
			\item Outsourcing the related process
			\item Insuring against a specific event
			\item Exiting the process that gives rise to risk
		\end{enumerate}
	\end{frame}
	\begin{frame}{D3-Q121: Answer}
		\begin{block}{Correct Answer: D}
			Avoidance means exiting the activities or conditions that give rise to risk. Examples include relocating a data center away from a hazard-prone region or declining to engage in a very large project with a notable risk of failure.
		\end{block}
	\end{frame}
	
	% ---- D3 Q122 ----
	\begin{frame}{D3-Q122: Addressing Critical Dependencies}
		\begin{block}{Question}
			Which would ensure that critical dependencies are addressed in the risk treatment plan?
		\end{block}
		\begin{enumerate}[(A)]
			\item Implement the risk treatment strategy for all possible risk
			\item Verify the accomplishment of business objectives through a top-down process review
			\item Treat each risk independently
			\item Verify the accomplishment of business objectives through a bottom-up process review
		\end{enumerate}
	\end{frame}
	\begin{frame}{D3-Q122: Answer}
		\begin{block}{Correct Answer: B}
			Development of an overall risk treatment strategy should be a top-down process, driven jointly by the need to achieve business objectives and to apply economically feasible constraints while controlling uncertainty within the enterprise risk appetite.
		\end{block}
	\end{frame}
	
	% ---- D3 Q123 ----
	\begin{frame}{D3-Q123: CISO Recommending Controls}
		\begin{block}{Question}
			A CISO recommends several controls, such as anti-malware, to protect the enterprise's information systems. Which approach to handling risk is being recommended?
		\end{block}
		\begin{enumerate}[(A)]
			\item Risk transference
			\item Risk mitigation
			\item Risk acceptance
			\item Risk avoidance
		\end{enumerate}
	\end{frame}
	\begin{frame}{D3-Q123: Answer}
		\begin{block}{Correct Answer: B --- Risk mitigation}
			By implementing controls, the company is trying to reduce the risk to an acceptable level, thereby mitigating risk.
		\end{block}
	\end{frame}
	
	% ---- D3 Q124 ----
	\begin{frame}{D3-Q124: Best Practice for Outsourcing Risk}
		\begin{block}{Question}
			Which practice BEST mitigates the risk associated with outsourcing a business function?
		\end{block}
		\begin{enumerate}[(A)]
			\item Performing audits to verify compliance with contract requirements
			\item Requiring all vendor staff to complete annual awareness training
			\item Retaining copies of all sensitive data on internal systems
			\item Reviewing the financial records of the vendor
		\end{enumerate}
	\end{frame}
	\begin{frame}{D3-Q124: Answer}
		\begin{block}{Correct Answer: A}
			When an outsourcing relationship is established, the risk of noncompliance with the agreement must be met through review, monitoring, and enforcement of contract terms. Regular audits verify vendor compliance.
		\end{block}
	\end{frame}
	
	% ---- D3 Q125 ----
	\begin{frame}{D3-Q125: Best Addressed by Transferring Risk}
		\begin{block}{Question}
			Which is BEST addressed by transferring risk?
		\end{block}
		\begin{enumerate}[(A)]
			\item An antiquated fire suppression system in the computer room
			\item The threat of disgruntled employee sabotage
			\item The threat of disgruntled employee theft
			\item A building located in a 100-year flood plain
		\end{enumerate}
	\end{frame}
	\begin{frame}{D3-Q125: Answer}
		\begin{block}{Correct Answer: D}
			Purchasing an insurance policy transfers the risk of a flood. Risk transfer is the process of assigning risk to another entity, usually through the purchase of an insurance policy or through outsourcing the service.
		\end{block}
	\end{frame}
	
	% ---- D3 Q126 ----
	\begin{frame}{D3-Q126: Finalizing Risk Treatment Plan}
		\begin{block}{Question}
			Which would BEST help finalize the risk treatment plan?
		\end{block}
		\begin{enumerate}[(A)]
			\item Vulnerability analysis
			\item Impact analysis
			\item Cost-benefit analysis
			\item SWOT analysis
		\end{enumerate}
	\end{frame}
	\begin{frame}{D3-Q126: Answer}
		\begin{block}{Correct Answer: C}
			A cost-benefit analysis helps determine if the benefit of a control outweighs the cost of implementing the control.
		\end{block}
	\end{frame}
	
	% ---- D3 Q127 ----
	\begin{frame}{D3-Q127: Greatest Benefit of Risk Treatment Plan}
		\begin{block}{Question}
			The GREATEST benefit of implementing a risk treatment plan is:
		\end{block}
		\begin{enumerate}[(A)]
			\item To reduce the impact and likelihood of risk occurrence
			\item To identify unmitigated risk that can be transferred
			\item To exploit the risk to test organizational preparedness
			\item To enhance the overall risk appetite of the enterprise
		\end{enumerate}
	\end{frame}
	\begin{frame}{D3-Q127: Answer}
		\begin{block}{Correct Answer: A}
			Implementing the risk treatment plan reduces the negative impact and likelihood of a risk occurrence.
		\end{block}
	\end{frame}
	
	% ---- D3 Q128 ----
	\begin{frame}{D3-Q128: Validating Patching Program}
		\begin{block}{Question}
			What is the BEST method to validate the effectiveness of an enterprise's patching program?
		\end{block}
		\begin{enumerate}[(A)]
			\item Conduct penetration testing
			\item Conduct a risk identification initiative
			\item Carry out vulnerability scans
			\item Review the requests for change
		\end{enumerate}
	\end{frame}
	\begin{frame}{D3-Q128: Answer}
		\begin{block}{Correct Answer: C}
			Performing vulnerability scans enables the IT risk practitioner to determine if patches are being installed on a timely basis. Penetration testing can elevate risk on production systems; RFCs do not prove patches are applied.
		\end{block}
	\end{frame}
	
	% ---- D3 Q129 ----
	\begin{frame}{D3-Q129: Designing IS Controls in Complex Environment}
		\begin{block}{Question}
			Which is the MOST important factor when designing information systems controls in a complex environment?
		\end{block}
		\begin{enumerate}[(A)]
			\item Development methodologies
			\item Scalability of the solution
			\item Technical platform interfaces
			\item Stakeholder requirements
		\end{enumerate}
	\end{frame}
	\begin{frame}{D3-Q129: Answer}
		\begin{block}{Correct Answer: D}
			The most important factor when designing IS controls is their ability to advance the interests of the business by addressing stakeholder requirements.
		\end{block}
	\end{frame}
	
	% ---- D3 Q130 ----
	\begin{frame}{D3-Q130: EFT Testing Focus}
		\begin{block}{Question}
			A financial institution is testing its EFT system after major enhancements. The risk practitioner would MOST benefit from a test that:
		\end{block}
		\begin{enumerate}[(A)]
			\item Identifies the introduction of potential new gaps in security
			\item Verifies adequate system recovery in case of failure
			\item Ensures the system performs to expectation
			\item Ensures the system can support transaction volume
		\end{enumerate}
	\end{frame}
	\begin{frame}{D3-Q130: Answer}
		\begin{block}{Correct Answer: A}
			The security of an EFT system is of the most interest to a risk practitioner. Testing should identify the introduction of potential new gaps in information security.
		\end{block}
	\end{frame}
	
	% ---- D3 Q131 ----
	\begin{frame}{D3-Q131: Creating KRIs for Quarterly Reporting}
		\begin{block}{Question}
			What is the BEST approach for creating KRIs for quarterly reporting to senior leadership?
		\end{block}
		\begin{enumerate}[(A)]
			\item Survey senior leaders about primary risk concerns
			\item Identify the most common network vulnerabilities
			\item Determine KRIs used in similar industry verticals
			\item Identify the enterprise risk appetite and metrics of current risk
		\end{enumerate}
	\end{frame}
	\begin{frame}{D3-Q131: Answer}
		\begin{block}{Correct Answer: D}
			KRIs are measurable and quantifiable indicators of risk. Creation of KRIs should be based on the enterprise's risk appetite and the measures available. Survey feedback is subjective; vulnerabilities lack likelihood/impact; other industries differ.
		\end{block}
	\end{frame}
	
	% ---- D3 Q132 ----
	\begin{frame}{D3-Q132: Gauging Stakeholder Concern}
		\begin{block}{Question}
			An enterprise contracted an external supplier for critical components. Risk tolerance levels for the outsourced component are documented and approved. Which can serve to gauge risk that may trigger stakeholder concern?
		\end{block}
		\begin{enumerate}[(A)]
			\item Indicators with approved thresholds
			\item Approved status report of completed milestones
			\item List of controls to be implemented by the supplier
			\item Number of findings in external audit reports
		\end{enumerate}
	\end{frame}
	\begin{frame}{D3-Q132: Answer}
		\begin{block}{Correct Answer: A}
			Indicators with approved thresholds demonstrate the acceptable risk levels stakeholders are willing to tolerate, and any risk above those approved levels will likely trigger stakeholder concern.
		\end{block}
	\end{frame}
	
	% ---- D3 Q133 ----
	\begin{frame}{D3-Q133: Revoking Temporary User Access}
		\begin{block}{Question}
			Which control within the user provision process BEST ensures revocation of system access for contractors and temporary users when no longer required?
		\end{block}
		\begin{enumerate}[(A)]
			\item Log account usage and send it to managers
			\item Establish predetermined, automatic expiration dates
			\item Ensure each user signs a security acknowledgment
			\item Require managers to email security when the user leaves
		\end{enumerate}
	\end{frame}
	\begin{frame}{D3-Q133: Answer}
		\begin{block}{Correct Answer: B}
			Predetermined expiration dates are the most effective means of removing systems access for temporary users. Logging is detective; acknowledgment has little effect; managers often delay termination notices.
		\end{block}
	\end{frame}
	
	% ---- D3 Q134 ----
	\begin{frame}{D3-Q134: Presenting Risk Profile to Board}
		\begin{block}{Question}
			Which BEST helps while presenting the current risk profile to executive management and the board?
		\end{block}
		\begin{enumerate}[(A)]
			\item Risk response dashboard
			\item Emerging risk report
			\item Risk register dashboard
			\item Key risk indicators report
		\end{enumerate}
	\end{frame}
	\begin{frame}{D3-Q134: Answer}
		\begin{block}{Correct Answer: C}
			A risk register dashboard provides a comprehensive overview of the risk profile of the enterprise. Risk response is a component; KRI report is only one component; emerging risk is not the full profile.
		\end{block}
	\end{frame}
	
	% ---- D3 Q135 ----
	\begin{frame}{D3-Q135: Senior Management Monitoring KCIs}
		\begin{block}{Question}
			What is the MAIN reason senior management monitors and analyzes trends in key control indicators?
		\end{block}
		\begin{enumerate}[(A)]
			\item It provides feedback on the overall control environment
			\item It helps in identifying redundant controls
			\item It proactively identifies impacts to the risk profile
			\item It helps determine if additional controls are required
		\end{enumerate}
	\end{frame}
	\begin{frame}{D3-Q135: Answer}
		\begin{block}{Correct Answer: C}
			The primary objective of KCIs is to ensure that controls actually mitigate risk at an effective level. Analysis of KCI trends provides information on the overall effectiveness of controls and the status of risk management.
		\end{block}
	\end{frame}
	
	% ---- D3 Q136 ----
	\begin{frame}{D3-Q136: Risk Treatment Plans Purpose}
		\begin{block}{Question}
			Risk treatment plans are necessary to describe how the:
		\end{block}
		\begin{enumerate}[(A)]
			\item Identified risk is further analyzed
			\item Chosen treatment options will be implemented
			\item Accepted risk is treated
			\item Risk indicators will monitor the risk
		\end{enumerate}
	\end{frame}
	\begin{frame}{D3-Q136: Answer}
		\begin{block}{Correct Answer: B}
			A risk treatment plan includes the plan of action for the chosen treatment, how it will be implemented, who will implement it, key dates, and resource requirements. Once risk is accepted, a treatment plan is not needed.
		\end{block}
	\end{frame}
	
	% ---- D3 Q137 ----
	\begin{frame}{D3-Q137: Cost-Benefit Analysis Stage}
		\begin{block}{Question}
			In the risk management process, a cost-benefit analysis is MAINLY performed:
		\end{block}
		\begin{enumerate}[(A)]
			\item As part of an initial risk assessment
			\item As part of risk-response planning
			\item During an information asset valuation
			\item When insurance is calculated for risk transfer
		\end{enumerate}
	\end{frame}
	\begin{frame}{D3-Q137: Answer}
		\begin{block}{Correct Answer: B}
			In risk response, a range of controls that can mitigate risk will be identified; a cost-benefit analysis helps identify the right controls to address the risk at acceptable levels within the budget.
		\end{block}
	\end{frame}
	
	% ---- D3 Q138 ----
	\begin{frame}{D3-Q138: Insufficient Resources for Treatment}
		\begin{block}{Question}
			In the event that available resources for risk treatment are not sufficient, the risk treatment plan should:
		\end{block}
		\begin{enumerate}[(A)]
			\item Define the priorities across all treatments to assist in resource allocation
			\item Recommend postponing treatment until resources are available
			\item Recommend reassessing the risk treatment options
			\item Suggest increasing the priority of the treatment
		\end{enumerate}
	\end{frame}
	\begin{frame}{D3-Q138: Answer}
		\begin{block}{Correct Answer: A}
			When available resources for treating the risk are not sufficient, the risk treatment plan should include prioritization of the risk so resources can be properly allocated.
		\end{block}
	\end{frame}
	
	% ---- D3 Q139 ----
	\begin{frame}{D3-Q139: System Backup and Restore Classification}
		\begin{block}{Question}
			System backup and restore procedures can BEST be classified as:
		\end{block}
		\begin{enumerate}[(A)]
			\item Technical controls
			\item Detective controls
			\item Corrective controls
			\item Deterrent controls
		\end{enumerate}
	\end{frame}
	\begin{frame}{D3-Q139: Answer}
		\begin{block}{Correct Answer: C --- Corrective controls}
			Corrective controls remediate the impact from negative events. If a system suffers harm so extensive that processing cannot continue, backup restore procedures enable that system to be recovered.
		\end{block}
	\end{frame}
	
	% ---- D3 Q140 ----
	\begin{frame}{D3-Q140: Primary Reason for Policy Exception}
		\begin{block}{Question}
			A PRIMARY reason for initiating a policy-exception process is:
		\end{block}
		\begin{enumerate}[(A)]
			\item The risk is justified by the benefit
			\item Policy compliance is difficult to enforce
			\item Operations are too busy to comply
			\item Users may initially be inconvenienced
		\end{enumerate}
	\end{frame}
	\begin{frame}{D3-Q140: Answer}
		\begin{block}{Correct Answer: A}
			Exceptions to policies are warranted if the benefits outweigh the costs of policy compliance; however, the enterprise needs to assess both the tangible and intangible risk and evaluate both in the context of existing risk.
		\end{block}
	\end{frame}
	
	% ---- D3 Q141 ----
	\begin{frame}{D3-Q141: Measuring Operational Controls}
		\begin{block}{Question}
			Which would BEST measure the effectiveness of operational controls?
		\end{block}
		\begin{enumerate}[(A)]
			\item Control matrix
			\item Key performance indicator
			\item Statement of applicability
			\item Key control indicator
		\end{enumerate}
	\end{frame}
	\begin{frame}{D3-Q141: Answer}
		\begin{block}{Correct Answer: D --- KCI}
			Key control indicators, also referred to as control effectiveness indicators, are metrics that provide information on the degree to which a control is working.
		\end{block}
	\end{frame}
	
	% ---- D3 Q142 ----
	\begin{frame}{D3-Q142: Selecting Security Technologies}
		\begin{block}{Question}
			Security technologies should be selected PRIMARILY on the basis of their:
		\end{block}
		\begin{enumerate}[(A)]
			\item Evaluation in security publications
			\item Compliance with industry standards
			\item Ability to mitigate risk to organizational objectives
			\item Cost compared to the enterprise's IT budget
		\end{enumerate}
	\end{frame}
	\begin{frame}{D3-Q142: Answer}
		\begin{block}{Correct Answer: C}
			The most fundamental criterion for selecting security technology is the capacity to reduce risk to organizational objectives.
		\end{block}
	\end{frame}
	
	% ---- D3 Q143 ----
	\begin{frame}{D3-Q143: Before Black Box Penetration Test}
		\begin{block}{Question}
			Which should be in place before a black box penetration test begins?
		\end{block}
		\begin{enumerate}[(A)]
			\item A clearly stated definition of scope
			\item Previous test results
			\item Proper communication and awareness training
			\item An incident response plan
		\end{enumerate}
	\end{frame}
	\begin{frame}{D3-Q143: Answer}
		\begin{block}{Correct Answer: A}
			A clearly stated definition of scope ensures a proper understanding of risk and success criteria. Previous results may inform scope; training and IR plan are not prerequisites.
		\end{block}
	\end{frame}
	
	% ---- D3 Q144 ----
	\begin{frame}{D3-Q144: Securing Corporate Network Against External Attack}
		\begin{block}{Question}
			What is the BEST way to ensure that a corporate network is adequately secured against external attack?
		\end{block}
		\begin{enumerate}[(A)]
			\item Use an intrusion detection system
			\item Establish minimum security baselines
			\item Implement vendor recommended settings
			\item Perform periodic penetration testing
		\end{enumerate}
	\end{frame}
	\begin{frame}{D3-Q144: Answer}
		\begin{block}{Correct Answer: D}
			Penetration testing is the best way to ensure that perimeter security is adequate. IDS is detective; baselines and vendor settings provide good practice but not assurance.
		\end{block}
	\end{frame}
	
	% ---- D3 Q145 ----
	\begin{frame}{D3-Q145: Before Establishing KCIs}
		\begin{block}{Question}
			Which is PRIMARILY defined before establishing key control indicators?
		\end{block}
		\begin{enumerate}[(A)]
			\item Desired capacity
			\item Desired risk threshold
			\item Desired tolerances
			\item Desired risk appetite
		\end{enumerate}
	\end{frame}
	\begin{frame}{D3-Q145: Answer}
		\begin{block}{Correct Answer: C --- Desired tolerances}
			The goal of KCIs is to track performance of control actions relative to tolerances, providing insight into the ongoing adequacy of a control in keeping risk within acceptable levels. Tolerances must be defined first.
		\end{block}
	\end{frame}
	
	% ---- D3 Q146 ----
	\begin{frame}{D3-Q146: Substantive Test of Tape Inventory}
		\begin{block}{Question}
			A substantive test to verify that tape library inventory records are accurate involves:
		\end{block}
		\begin{enumerate}[(A)]
			\item Determining whether bar code readers are installed
			\item Conducting a physical count of the tape inventory
			\item Checking whether receipts and issues of tapes are accurately recorded
			\item Determining whether movement of tapes is authorized
		\end{enumerate}
	\end{frame}
	\begin{frame}{D3-Q146: Answer}
		\begin{block}{Correct Answer: B}
			A substantive test involves gathering evidence to evaluate the integrity of individual transactions, data, or other information. A physical count of the tape inventory is a substantive test; the others are compliance tests.
		\end{block}
	\end{frame}
	
	% ---- D3 Q147 ----
	\begin{frame}{D3-Q147: Ensuring Mitigation of Vulnerabilities}
		\begin{block}{Question}
			Which BEST ensures that appropriate mitigation occurs on identified information systems vulnerabilities?
		\end{block}
		\begin{enumerate}[(A)]
			\item Presenting root cause analysis to management
			\item Implementing software to input action points
			\item Incorporating findings into the annual report
			\item Assigning action plans with deadlines to responsible personnel
		\end{enumerate}
	\end{frame}
	\begin{frame}{D3-Q147: Answer}
		\begin{block}{Correct Answer: D}
			Assigning mitigation to personnel establishes responsibility for its completion within the deadline. Presenting findings raises awareness but does not ensure action; software and reporting do not guarantee completion.
		\end{block}
	\end{frame}
	
	% ---- D3 Q148 ----
	\begin{frame}{D3-Q148: Audit Report on Uncritical Functions}
		\begin{block}{Question}
			What is the BEST action if an audit report recommends supplemental controls for uncritical functions?
		\end{block}
		\begin{enumerate}[(A)]
			\item Appoint a project lead to start implementation
			\item Eliminate the risk by purchasing risk insurance
			\item Request that IT implement all supplemental controls
			\item Request a cost-benefit analysis to determine business need
		\end{enumerate}
	\end{frame}
	\begin{frame}{D3-Q148: Answer}
		\begin{block}{Correct Answer: D}
			A cost-benefit analysis may help determine if supplemental controls are needed, ensuring the cost of the control does not exceed the impact from the risk being realized.
		\end{block}
	\end{frame}
	
	% ---- D3 Q149 ----
	\begin{frame}{D3-Q149: Peer Review of Risk Management}
		\begin{block}{Question}
			Which BEST enables a peer review of an enterprise's risk management process?
		\end{block}
		\begin{enumerate}[(A)]
			\item A balanced scorecard
			\item An industry survey
			\item A capability maturity model
			\item A framework
		\end{enumerate}
	\end{frame}
	\begin{frame}{D3-Q149: Answer}
		\begin{block}{Correct Answer: C}
			A capability maturity model describes essential elements and criteria for effective processes and outlines an evolutionary improvement path from ad hoc to disciplined, mature processes.
		\end{block}
	\end{frame}
	
	% ---- D3 Q150 ----
	\begin{frame}{D3-Q150: Risk Treatment Plan With Resolution}
		\begin{block}{Question}
			Which is the MOST important information to include in a risk treatment plan that already has an appropriate resolution and a date for completion?
		\end{block}
		\begin{enumerate}[(A)]
			\item Responsible personnel
			\item Mitigating factors
			\item Likelihood of occurrence
			\item Cost of completion
		\end{enumerate}
	\end{frame}
	\begin{frame}{D3-Q150: Answer}
		\begin{block}{Correct Answer: A}
			Risk response activities must be assigned to a responsible person or group; without this assignment, it will be unclear who will implement the countermeasure.
		\end{block}
	\end{frame}
	
	% ---- D3 Q151 ----
	\begin{frame}{D3-Q151: Ensuring Risk Remains Acceptable}
		\begin{block}{Question}
			Which BEST ensures that identified risk remains at an acceptable level?
		\end{block}
		\begin{enumerate}[(A)]
			\item Reviewing controls periodically according to the risk treatment plan
			\item Listing each risk as a separate entry in the risk register
			\item Creating a separate risk register for every department
			\item Maintaining a KRI for assets in the risk register
		\end{enumerate}
	\end{frame}
	\begin{frame}{D3-Q151: Answer}
		\begin{block}{Correct Answer: A}
			Controls deployed according to the risk treatment plan should provide the desired results, because the plan is based on management's acceptance of residual risk and approval of deployment steps.
		\end{block}
	\end{frame}
	
	% ---- D3 Q152 ----
	\begin{frame}{D3-Q152: Risk Response Increasing Liability}
		\begin{block}{Question}
			Which risk response option is MOST likely to increase the liability of the enterprise?
		\end{block}
		\begin{enumerate}[(A)]
			\item Risk acceptance
			\item Risk reduction
			\item Risk transfer
			\item Risk avoidance
		\end{enumerate}
	\end{frame}
	\begin{frame}{D3-Q152: Answer}
		\begin{block}{Correct Answer: A --- Risk acceptance}
			An enterprise may choose to accept risk without knowing the correct level of risk being accepted; this may result in accusations of negligence.
		\end{block}
	\end{frame}
	
	% ---- D3 Q153 ----
	\begin{frame}{D3-Q153: Best Reason Not to Reduce Risk}
		\begin{block}{Question}
			Which is the BEST reason for an enterprise to decide not to reduce an identified risk?
		\end{block}
		\begin{enumerate}[(A)]
			\item There is no regulatory requirement to reduce the risk
			\item There are mitigating controls in place
			\item The cost of mitigation exceeds the risk
			\item The budget for risk mitigation is limited
		\end{enumerate}
	\end{frame}
	\begin{frame}{D3-Q153: Answer}
		\begin{block}{Correct Answer: C}
			Enterprises will accept the risk when the cost of mitigation exceeds the risk. Regulatory requirement, existing controls (residual may still exceed tolerance), and budget alone are insufficient reasons.
		\end{block}
	\end{frame}
	
	% ---- D3 Q154 ----
	\begin{frame}{D3-Q154: Outsourcing Part of IT Activities}
		\begin{block}{Question}
			An enterprise decides to address risk associated with an IT project by outsourcing part of the IT activities to a third party with a specialized skill set. In relation to the project, this is an example of:
		\end{block}
		\begin{enumerate}[(A)]
			\item Risk transfer
			\item Risk avoidance
			\item Risk acceptance
			\item Risk mitigation
		\end{enumerate}
	\end{frame}
	\begin{frame}{D3-Q154: Answer}
		\begin{block}{Correct Answer: D --- Risk mitigation}
			When specific activities are outsourced to an entity with a specialized skill set, the inherent risk of the activity is reduced. The outsourcing process works as a control in this case.
		\end{block}
	\end{frame}
	
	% ---- D3 Q155 ----
	\begin{frame}{D3-Q155: Facilitating Cost-Effective Risk Response}
		\begin{block}{Question}
			Which BEST facilitates cost-effective risk response?
		\end{block}
		\begin{enumerate}[(A)]
			\item Prioritizing and addressing risk according to the risk management strategy
			\item Mitigating risk on the basis of likelihood and magnitude of impact
			\item Performing countermeasure analysis for each control deployed
			\item Selecting controls that are at zero or near-zero costs
		\end{enumerate}
	\end{frame}
	\begin{frame}{D3-Q155: Answer}
		\begin{block}{Correct Answer: A}
			Prioritizing and addressing risk in line with the risk treatment strategy balances the costs and benefits of managing the IT risk. Likelihood/impact alone do not optimize treatment; zero-cost controls may be ineffective.
		\end{block}
	\end{frame}
	
	% ---- D3 Q156 ----
	\begin{frame}{D3-Q156: Recommending DLP Device}
		\begin{block}{Question}
			Which can BEST be used as a basis for recommending a data loss prevention (DLP) device as a security control?
		\end{block}
		\begin{enumerate}[(A)]
			\item Benchmarking with peers on DLP deployment
			\item A business case for DLP to protect data
			\item Evaluation report of popular DLP solutions
			\item Data leakage scenario in risk register
		\end{enumerate}
	\end{frame}
	\begin{frame}{D3-Q156: Answer}
		\begin{block}{Correct Answer: B}
			A business case with costs versus benefits provides the business reasoning for why the DLP solution addresses the risk and explains how risk losses could be reduced.
		\end{block}
	\end{frame}
	
	% ---- D3 Q157 ----
	\begin{frame}{D3-Q157: Cost Greater Than Opportunity}
		\begin{block}{Question}
			When the cost of risk related to a specific business process is greater than the potential opportunity, the BEST risk response is:
		\end{block}
		\begin{enumerate}[(A)]
			\item Transfer
			\item Acceptance
			\item Mitigation
			\item Avoidance
		\end{enumerate}
	\end{frame}
	\begin{frame}{D3-Q157: Answer}
		\begin{block}{Correct Answer: D --- Avoidance}
			Risk avoidance is the process for systematically avoiding risk, constituting one approach to managing risk. When risk cost exceeds the potential opportunity, avoiding the activity is best.
		\end{block}
	\end{frame}
	
	% ---- D3 Q158 ----
	\begin{frame}{D3-Q158: First Step When Outsourcing All IT}
		\begin{block}{Question}
			What is the FIRST step for a risk practitioner when an enterprise decides to outsource all IT services and support to a third party?
		\end{block}
		\begin{enumerate}[(A)]
			\item Validate that the internal systems of the provider are secure
			\item Enforce transborder data flow regulations
			\item Ensure security requirements are addressed in all contracts and agreements
			\item Build a business case to perform an on-site audit of the vendor
		\end{enumerate}
	\end{frame}
	\begin{frame}{D3-Q158: Answer}
		\begin{block}{Correct Answer: C}
			A contract only covers topics listed in the contract. If security is not explicitly included in the contract terms, the enterprise may not be properly protected.
		\end{block}
	\end{frame}
	
	% ---- D3 Q159 ----
	\begin{frame}{D3-Q159: Justifying IT Control Cost}
		\begin{block}{Question}
			Which BEST helps ensure that the cost is justifiable when selecting an IT control?
		\end{block}
		\begin{enumerate}[(A)]
			\item The investment cost is within budget
			\item The risk likelihood and impact are reflected
			\item The net present value of the IT control cost is high
			\item Low-cost open source technology is used
		\end{enumerate}
	\end{frame}
	\begin{frame}{D3-Q159: Answer}
		\begin{block}{Correct Answer: B}
			While other factors may be relevant, the total cost of ownership of a control should not exceed the projected likelihood times the impact of the risk it is intended to mitigate.
		\end{block}
	\end{frame}
	
	% ---- D3 Q160 ----
	\begin{frame}{D3-Q160: Testing Compliance of Response Plan}
		\begin{block}{Question}
			Testing compliance of a response and recovery plan should begin by conducting a:
		\end{block}
		\begin{enumerate}[(A)]
			\item Tabletop exercise
			\item Review of archived logs
			\item Penetration test
			\item Business impact analysis
		\end{enumerate}
	\end{frame}
	\begin{frame}{D3-Q160: Answer}
		\begin{block}{Correct Answer: A}
			Tabletop exercises simulate incidents to test the response capability of an enterprise. Participants gather to formulate responses to each development in realistic scenarios.
		\end{block}
	\end{frame}
	
	% ---- D3 Q161 ----
	\begin{frame}{D3-Q161: Capability Maturity Model Basis}
		\begin{block}{Question}
			The use of a capability maturity model is based on:
		\end{block}
		\begin{enumerate}[(A)]
			\item The training of staff
			\item The development of new controls
			\item The application of standard, repeatable processes that can be measured
			\item Users developing new innovative solutions
		\end{enumerate}
	\end{frame}
	\begin{frame}{D3-Q161: Answer}
		\begin{block}{Correct Answer: C}
			The use of maturity models requires development, implementation, and measuring of consistent procedures and activities. Maturity models describe an evolutionary path from ad hoc to mature, measured processes.
		\end{block}
	\end{frame}
	
	% ---- D3 Q162 ----
	\begin{frame}{D3-Q162: Internal Control Requirements Source}
		\begin{block}{Question}
			An enterprise is implementing controls to protect employee details from unauthorized exposure. The internal control requirements will come from:
		\end{block}
		\begin{enumerate}[(A)]
			\item The risk management team
			\item Internal audit
			\item IT management
			\item Process owners
		\end{enumerate}
	\end{frame}
	\begin{frame}{D3-Q162: Answer}
		\begin{block}{Correct Answer: D --- Process owners}
			Process owners provide the internal control requirements based on the business needs and objectives. Risk management reviews control risk; audit measures effectiveness; IT handles IT governance.
		\end{block}
	\end{frame}
	
	% ---- D3 Q163 ----
	\begin{frame}{D3-Q163: Costs Balanced with Potential Loss}
		\begin{block}{Question}
			Which costs MUST be balanced with the potential for loss to ensure effective risk management?
		\end{block}
		\begin{enumerate}[(A)]
			\item Regulatory fines and penalties
			\item Total cost of ownership
			\item Breach management cost
			\item Business operation cost
		\end{enumerate}
	\end{frame}
	\begin{frame}{D3-Q163: Answer}
		\begin{block}{Correct Answer: B}
			Total cost of ownership considers implementation and operating costs for a control, which must be balanced with the potential for loss to ensure effective risk management.
		\end{block}
	\end{frame}
	
	% ---- D3 Q164 ----
	\begin{frame}{D3-Q164: Controls Design Reduction}
		\begin{block}{Question}
			Controls are most effective when they are designed to reduce:
		\end{block}
		\begin{enumerate}[(A)]
			\item Threats
			\item Likelihood
			\item Uncertainty
			\item Vulnerabilities
		\end{enumerate}
	\end{frame}
	\begin{frame}{D3-Q164: Answer}
		\begin{block}{Correct Answer: D --- Vulnerabilities}
			Controls are most effective when designed to reduce vulnerabilities affecting the enterprise. Threats cannot always be reduced; likelihood alone is one aspect; controls generally do not reduce uncertainty.
		\end{block}
	\end{frame}
	
	% ---- D3 Q165 ----
	\begin{frame}{D3-Q165: Mitigating Patch Interoperability Risk}
		\begin{block}{Question}
			System administrators may release critical patches into production without testing. Which would BEST mitigate the risk of interoperability issues?
		\end{block}
		\begin{enumerate}[(A)]
			\item Ensure a reliable system rollback plan is in place
			\item Test the patch on the least critical systems first
			\item Only allow updates to occur after hours
			\item Ensure patches are approved by the CISO
		\end{enumerate}
	\end{frame}
	\begin{frame}{D3-Q165: Answer}
		\begin{block}{Correct Answer: A}
			A reliable system rollback plan allows administrators to roll back patches in case they affect the system negatively. Testing on least critical systems does not reveal impact on all applications.
		\end{block}
	\end{frame}
	
	% ============================================================
	% DOMAIN 4 --- INFORMATION TECHNOLOGY AND SECURITY (22%) --- Q166--Q200
	% ============================================================
	\section{Domain 4: Information Technology and Security}
	
	% ---- D4 Q166 ----
	\begin{frame}{D4-Q166: Information No Longer Required}
		\begin{block}{Question}
			Information no longer required to support the main purpose of the enterprise should be managed:
		\end{block}
		\begin{enumerate}[(A)]
			\item Under the retention policy
			\item Under the information classification policy
			\item Under the backup policy
			\item Under the business impact analysis
		\end{enumerate}
	\end{frame}
	\begin{frame}{D4-Q166: Answer}
		\begin{block}{Correct Answer: A --- Retention policy}
			Information no longer required should be analyzed under the retention policy to determine whether the enterprise is required to maintain the data for business, legal, or regulatory reasons. Keeping data unnecessarily increases risk.
		\end{block}
	\end{frame}
	
	% ---- D4 Q167 ----
	\begin{frame}{D4-Q167: Protecting Breakthrough Technology}
		\begin{block}{Question}
			An enterprise developed a breakthrough technology providing a competitive edge. Which FIRST governs how this information is protected from within the enterprise?
		\end{block}
		\begin{enumerate}[(A)]
			\item The data classification policy
			\item The acceptable use policy
			\item Encryption standards
			\item The access control policy
		\end{enumerate}
	\end{frame}
	\begin{frame}{D4-Q167: Answer}
		\begin{block}{Correct Answer: A}
			A data classification policy describes the data classification categories, the level of protection for each category, and roles and responsibilities. Encryption and access controls are driven by the classification.
		\end{block}
	\end{frame}
	
	% ---- D4 Q168 ----
	\begin{frame}{D4-Q168: Most Important IS Infrastructure Requirement}
		\begin{block}{Question}
			Which is the MOST important requirement for setting up an information security infrastructure for a new system?
		\end{block}
		\begin{enumerate}[(A)]
			\item Performing a business impact analysis
			\item Considering personal devices as part of the security policy
			\item Basing the information security infrastructure on a risk assessment
			\item Initiating an IT security awareness campaign
		\end{enumerate}
	\end{frame}
	\begin{frame}{D4-Q168: Answer}
		\begin{block}{Correct Answer: C}
			The information security infrastructure should be based on a risk assessment. BIA is for continuity; personal devices may be part of policy; awareness training is supportive.
		\end{block}
	\end{frame}
	
	% ---- D4 Q169 ----
	\begin{frame}{D4-Q169: Data Retention Policy Primary Concern}
		\begin{block}{Question}
			The PRIMARY concern of a risk practitioner reviewing a formal data retention policy is:
		\end{block}
		\begin{enumerate}[(A)]
			\item Storage availability
			\item Applicable organizational standards
			\item Generally accepted industry good practices
			\item Regulatory and business requirements
		\end{enumerate}
	\end{frame}
	\begin{frame}{D4-Q169: Answer}
		\begin{block}{Correct Answer: D}
			In determining the retention policy, regulatory requirements are of primary importance along with business requirements. Without them, an enterprise can keep records indefinitely at tremendous cost.
		\end{block}
	\end{frame}
	
	% ---- D4 Q170 ----
	\begin{frame}{D4-Q170: Security Awareness Programs Reason}
		\begin{block}{Question}
			Which is the MOST important reason for conducting security awareness programs throughout an enterprise?
		\end{block}
		\begin{enumerate}[(A)]
			\item Reducing the risk of social engineering attacks
			\item Training personnel to respond to security incidents
			\item Informing business units about the security strategy
			\item Maintaining evidence of training records
		\end{enumerate}
	\end{frame}
	\begin{frame}{D4-Q170: Answer}
		\begin{block}{Correct Answer: A}
			Security awareness reduces the risk of successful social engineering attacks by sensitizing employees to security policies and risk. People are often the weakest link in security implementations.
		\end{block}
	\end{frame}
	
	% ---- D4 Q171 ----
	\begin{frame}{D4-Q171: Incident Response Training Effective}
		\begin{block}{Question}
			Which is the BEST indicator that incident response training is effective?
		\end{block}
		\begin{enumerate}[(A)]
			\item Decreased reporting of security incidents
			\item Increased reporting of security incidents
			\item Decreased number of password resets
			\item Increased number of identified system vulnerabilities
		\end{enumerate}
	\end{frame}
	\begin{frame}{D4-Q171: Answer}
		\begin{block}{Correct Answer: B}
			Increased reporting of valid incidents is the best indicator because it shows users are aware of security rules and know how to report incidents. Decreased reporting suggests lack of awareness.
		\end{block}
	\end{frame}
	
	% ---- D4 Q172 ----
	\begin{frame}{D4-Q172: Greatest Information Security Awareness}
		\begin{block}{Question}
			Which provides the GREATEST level of information security awareness?
		\end{block}
		\begin{enumerate}[(A)]
			\item Job descriptions
			\item A security manual
			\item Security training
			\item An organizational diagram
		\end{enumerate}
	\end{frame}
	\begin{frame}{D4-Q172: Answer}
		\begin{block}{Correct Answer: C}
			Security training is the best way to inform all employees about changes to the risk landscape and enhance awareness of enterprise risk management strategy.
		\end{block}
	\end{frame}
	
	% ---- D4 Q173 ----
	\begin{frame}{D4-Q173: Shared Generic Account}
		\begin{block}{Question}
			Which is the MAIN concern when two or more staff members are allowed to use the same generic account?
		\end{block}
		\begin{enumerate}[(A)]
			\item Segregation of duties
			\item Inability to change the password
			\item Repudiation
			\item Inability to trace account activities
		\end{enumerate}
	\end{frame}
	\begin{frame}{D4-Q173: Answer}
		\begin{block}{Correct Answer: C --- Repudiation}
			Because username and password are the same for generic accounts, repudiation becomes an issue. None of the users can be held accountable because each can deny accountability for transactions.
		\end{block}
	\end{frame}
	
	% ---- D4 Q174 ----
	\begin{frame}{D4-Q174: IT Risk Awareness Campaign Approach}
		\begin{block}{Question}
			Which is the BEST approach when conducting an IT risk awareness campaign?
		\end{block}
		\begin{enumerate}[(A)]
			\item Provide technical detail on exploits
			\item Provide common messages tailored for different groups
			\item Target system administrators and help desk staff
			\item Target senior managers and business process owners
		\end{enumerate}
	\end{frame}
	\begin{frame}{D4-Q174: Answer}
		\begin{block}{Correct Answer: B}
			Groups differ in level of responsibility and expertise; tailor common messages to each group's role and understanding. Providing exploit details could enable circumvention; targeting specific groups excludes others.
		\end{block}
	\end{frame}
	
	% ---- D4 Q175 ----
	\begin{frame}{D4-Q175: Social Media Outlet Greatest Concern}
		\begin{block}{Question}
			Which principle of information security is of GREATEST concern to a social media outlet for political or news-related content?
		\end{block}
		\begin{enumerate}[(A)]
			\item Integrity
			\item Confidentiality
			\item Availability
			\item Nonrepudiation
		\end{enumerate}
	\end{frame}
	\begin{frame}{D4-Q175: Answer}
		\begin{block}{Correct Answer: C --- Availability}
			For a social media outlet, availability is of the greatest concern. Integrity, confidentiality, and nonrepudiation are typically the responsibility of social media users.
		\end{block}
	\end{frame}
	
	% ---- D4 Q176 ----
	\begin{frame}{D4-Q176: Most Significant Threat to Project}
		\begin{block}{Question}
			Which poses the MOST significant threat to a project?
		\end{block}
		\begin{enumerate}[(A)]
			\item A lack of feedback upon project completion
			\item A lack of unit testing
			\item Missed opportunities from lessons learned
			\item Misunderstanding the requirements
		\end{enumerate}
	\end{frame}
	\begin{frame}{D4-Q176: Answer}
		\begin{block}{Correct Answer: D}
			Projects exist to deliver specific outcomes. If requirements are misunderstood, a project can succeed internally (schedule, budget) yet fail the business because it did not deliver business value.
		\end{block}
	\end{frame}
	
	% ---- D4 Q177 ----
	\begin{frame}{D4-Q177: SDLC Phase for Risk Practitioner}
		\begin{block}{Question}
			In which phase of the SDLC should a risk practitioner FIRST become involved?
		\end{block}
		\begin{enumerate}[(A)]
			\item Analysis
			\item Design
			\item Planning
			\item Implementation
		\end{enumerate}
	\end{frame}
	\begin{frame}{D4-Q177: Answer}
		\begin{block}{Correct Answer: C --- Planning}
			The risk practitioner should become involved as early as possible. When participating in planning, they can influence resource requests to meet risk objectives most efficiently and effectively.
		\end{block}
	\end{frame}
	
	% ---- D4 Q178 ----
	\begin{frame}{D4-Q178: Project Going Over Budget}
		\begin{block}{Question}
			Which is MOST likely to result in a project going over budget?
		\end{block}
		\begin{enumerate}[(A)]
			\item Testing
			\item Scope creep
			\item Loss of sponsorship
			\item Quality control
		\end{enumerate}
	\end{frame}
	\begin{frame}{D4-Q178: Answer}
		\begin{block}{Correct Answer: B --- Scope creep}
			Scope creep refers to uncontrolled changes in a project's scope. Unless scope is controlled, duration and budget cannot be held to account; the project may exceed budget to meet changing requirements.
		\end{block}
	\end{frame}
	
	% ---- D4 Q179 ----
	\begin{frame}{D4-Q179: Project Risk Identification}
		\begin{block}{Question}
			Project-related risk that can threaten the successful completion of a software development project is BEST identified by the:
		\end{block}
		\begin{enumerate}[(A)]
			\item Project manager
			\item Project sponsor
			\item Project steering committee
			\item Quality assurance team
		\end{enumerate}
	\end{frame}
	\begin{frame}{D4-Q179: Answer}
		\begin{block}{Correct Answer: A --- Project manager}
			The project manager is responsible for identification, analysis, and monitoring of project-related risk. Risk is reported to the project sponsor.
		\end{block}
	\end{frame}
	
	% ---- D4 Q180 ----
	\begin{frame}{D4-Q180: VPN Implementation Concern}
		\begin{block}{Question}
			Which is of MOST concern in a review of a VPN implementation?
		\end{block}
		\begin{enumerate}[(A)]
			\item Computers on the network are located at the enterprise's remote offices
			\item Computers are on the enterprise's internal network
			\item Computers are at the backup site
			\item Computers are in employees' homes
		\end{enumerate}
	\end{frame}
	\begin{frame}{D4-Q180: Answer}
		\begin{block}{Correct Answer: D}
			Home computers are least often subject to the enterprise security policy and therefore are high-risk. Once a computer is compromised, any network that trusts it is at risk.
		\end{block}
	\end{frame}
	
	% ---- D4 Q181 ----
	\begin{frame}{D4-Q181: BCM Meeting Stakeholder Expectations}
		\begin{block}{Question}
			Which BEST supports business continuity management in meeting external stakeholder expectations?
		\end{block}
		\begin{enumerate}[(A)]
			\item Prioritizing applications based on business criticality
			\item Ensuring backup data are available for restoration
			\item Disclosing the crisis management strategy statement
			\item Obtaining risk assessment by an independent party
		\end{enumerate}
	\end{frame}
	\begin{frame}{D4-Q181: Answer}
		\begin{block}{Correct Answer: A}
			External parties expect their information assets are secured. Prioritizing applications based on business criticality allows external expectations to be met optimally with limited resources.
		\end{block}
	\end{frame}
	
	% ---- D4 Q182 ----
	\begin{frame}{D4-Q182: Cloud Computing Risk Controls}
		\begin{block}{Question}
			Which BEST identifies controls addressing risk related to cloud computing?
		\end{block}
		\begin{enumerate}[(A)]
			\item Data encryption, tenant isolation, controlled change management
			\item Data encryption, customizing the application template
			\item Use of open standards technology, data encryption, tenant isolation
			\item Tenant isolation, controlled change management, custom widgets
		\end{enumerate}
	\end{frame}
	\begin{frame}{D4-Q182: Answer}
		\begin{block}{Correct Answer: A}
			Encryption facilitates separation of data among tenants. Tenant isolation ensures that one tenant's data are sequestered from other tenants. Controlled change management ensures all changes are planned and tenant dependencies are mapped.
		\end{block}
	\end{frame}
	
	% ---- D4 Q183 ----
	\begin{frame}{D4-Q183: Self-Signed Digital Certificates}
		\begin{block}{Question}
			A financial institution conducts high-value electronic transactions using customers' self-signed digital certificates. Which is of MOST concern?
		\end{block}
		\begin{enumerate}[(A)]
			\item Self-signed certificates may not be accepted by internal browsers
			\item Nonrepudiation requirements may not be enforced
			\item The certificate's private key may be susceptible to attack
			\item Self-signed certificates may lead to session hijacking
		\end{enumerate}
	\end{frame}
	\begin{frame}{D4-Q183: Answer}
		\begin{block}{Correct Answer: B}
			Nonrepudiation of transactions is critical in financial transactions. This can only be satisfied by using third party-trusted certificates signed by an accredited authority.
		\end{block}
	\end{frame}
	
	% ---- D4 Q184 ----
	\begin{frame}{D4-Q184: Hardware vs Software Risk}
		\begin{block}{Question}
			Which risk consideration is of MOST concern to the risk practitioner regarding hardware as opposed to software?
		\end{block}
		\begin{enumerate}[(A)]
			\item Misconfiguration
			\item Access control
			\item Skilled resources
			\item Version control
		\end{enumerate}
	\end{frame}
	\begin{frame}{D4-Q184: Answer}
		\begin{block}{Correct Answer: A --- Misconfiguration}
			Misconfiguration can lead to potential vulnerabilities, leaving the system open to external parties, and is a larger concern for hardware than software. Access control and skilled resources apply to both; version control does not apply to hardware.
		\end{block}
	\end{frame}
	
	% ---- D4 Q185 ----
	\begin{frame}{D4-Q185: System Requirements Risk Phase}
		\begin{block}{Question}
			At which phase of the SDLC should risk related to system requirements be determined?
		\end{block}
		\begin{enumerate}[(A)]
			\item Development
			\item Implementation
			\item Initiation
			\item Maintenance
		\end{enumerate}
	\end{frame}
	\begin{frame}{D4-Q185: Answer}
		\begin{block}{Correct Answer: C --- Initiation}
			Upon initiation of the SDLC, identified risk should inform system requirements, including security requirements and a security concept of operations.
		\end{block}
	\end{frame}
	
	% ---- D4 Q186 ----
	\begin{frame}{D4-Q186: Wrong Patient File}
		\begin{block}{Question}
			A healthcare practitioner providing care is given the file of a patient scheduled for an appointment later that day. What is the PRIMARY type of risk faced?
		\end{block}
		\begin{enumerate}[(A)]
			\item Relevance
			\item Integrity
			\item Security
			\item Availability
		\end{enumerate}
	\end{frame}
	\begin{frame}{D4-Q186: Answer}
		\begin{block}{Correct Answer: A --- Relevance risk}
			Relevance risk is the risk of not getting the right information to the right person at the right time for the correct action. The file provided is not relevant to that patient.
		\end{block}
	\end{frame}
	
	% ---- D4 Q187 ----
	\begin{frame}{D4-Q187: Best Cloud Model for Confidential Info}
		\begin{block}{Question}
			An enterprise accepts the risk of confidential information in the cloud. Which cloud deployment model offers the most safeguards?
		\end{block}
		\begin{enumerate}[(A)]
			\item Public cloud
			\item Private cloud
			\item Community cloud
			\item Hybrid cloud
		\end{enumerate}
	\end{frame}
	\begin{frame}{D4-Q187: Answer}
		\begin{block}{Correct Answer: B --- Private cloud}
			The private cloud model operates solely for the enterprise and will have controls in place needed to keep enterprise information confidential.
		\end{block}
	\end{frame}
	
	% ---- D4 Q188 ----
	\begin{frame}{D4-Q188: BCP Authorship}
		\begin{block}{Question}
			Business continuity plans should be written and maintained by:
		\end{block}
		\begin{enumerate}[(A)]
			\item The information security and IT functions
			\item Representatives from all functional units
			\item The risk management function
			\item Executive management
		\end{enumerate}
	\end{frame}
	\begin{frame}{D4-Q188: Answer}
		\begin{block}{Correct Answer: B}
			Business continuity planning is an enterprise-wide activity; it is successful only if all business owners collaborate in developing, testing, and maintaining the plan. IS/IT handle DR; risk management oversees; executives ensure planning.
		\end{block}
	\end{frame}
	
	% ---- D4 Q189 ----
	\begin{frame}{D4-Q189: Verifying Up-to-Date Antivirus}
		\begin{block}{Question}
			What is the BEST way to verify that critical production servers are using up-to-date antivirus signature files?
		\end{block}
		\begin{enumerate}[(A)]
			\item Check a sample of servers
			\item Verify the date signature files were last pushed out
			\item Use a benign virus to test quarantine
			\item Research the most recent signature file and compare to the console
		\end{enumerate}
	\end{frame}
	\begin{frame}{D4-Q189: Answer}
		\begin{block}{Correct Answer: A}
			The only effective way to verify currency of signature files is to look at a sample of servers. Push-out date does not guarantee loading; releasing a virus is inappropriate; console comparison does not prove loading.
		\end{block}
	\end{frame}
	
	% ---- D4 Q190 ----
	\begin{frame}{D4-Q190: Security Awareness Program Effectiveness}
		\begin{block}{Question}
			Which is MOST important for measuring the effectiveness of a security awareness program?
		\end{block}
		\begin{enumerate}[(A)]
			\item Increased interest in focus groups on security
			\item A reduced number of security violation reports
			\item A quantitative evaluation to ensure user comprehension
			\item An increased number of security violation reports
		\end{enumerate}
	\end{frame}
	\begin{frame}{D4-Q190: Answer}
		\begin{block}{Correct Answer: D}
			An increase in the number of violation reports is the best indicator of a high level of security awareness. As with automated alerts, each report needs to be assessed for validity.
		\end{block}
	\end{frame}
	
	% ---- D4 Q191 ----
	\begin{frame}{D4-Q191: Major Risk of Single Sign-On}
		\begin{block}{Question}
			A MAJOR risk of single sign-on is that it:
		\end{block}
		\begin{enumerate}[(A)]
			\item Uses complex technologies for password management
			\item May potentially bypass the enterprise firewall
			\item Is prone to distributed denial-of-service attacks
			\item May be a potential single point of compromise
		\end{enumerate}
	\end{frame}
	\begin{frame}{D4-Q191: Answer}
		\begin{block}{Correct Answer: D}
			Because SSO uses a single username and password for access to multiple systems, loss or disclosure of the credentials may compromise a broad range of systems.
		\end{block}
	\end{frame}
	
	% ---- D4 Q192 ----
	\begin{frame}{D4-Q192: Changeover Consistency of Data}
		\begin{block}{Question}
			Which changeover technique would have the GREATEST risk regarding consistency of data when moving users from one system to another?
		\end{block}
		\begin{enumerate}[(A)]
			\item Parallel
			\item Outsourcing
			\item Abrupt
			\item Phased
		\end{enumerate}
	\end{frame}
	\begin{frame}{D4-Q192: Answer}
		\begin{block}{Correct Answer: C --- Abrupt}
			Abrupt changeover occurs when users are converted immediately upon availability. It is usually least expensive but involves high risk of data loss and system failure, with no fallback to the old system.
		\end{block}
	\end{frame}
	
	% ---- D4 Q193 ----
	\begin{frame}{D4-Q193: Insider Threat to Confidential Info}
		\begin{block}{Question}
			Which is the MOST effective measure against insider threats to confidential information?
		\end{block}
		\begin{enumerate}[(A)]
			\item Audit trail monitoring
			\item A privacy policy
			\item Role-based access control
			\item Defense in depth
		\end{enumerate}
	\end{frame}
	\begin{frame}{D4-Q193: Answer}
		\begin{block}{Correct Answer: C}
			Role-based access control provides access according to business need, reducing unnecessary access rights and enforcing accountability, helping guard against breaches of confidential information.
		\end{block}
	\end{frame}
	
	% ---- D4 Q194 ----
	\begin{frame}{D4-Q194: Improving Awareness After Spear Phishing}
		\begin{block}{Question}
			Despite a comprehensive awareness program, a breach occurred via spear phishing. What is the MOST effective way to improve security awareness?
		\end{block}
		\begin{enumerate}[(A)]
			\item Review and improve coverage of social engineering threats
			\item Launch a disciplinary process against those who leaked information
			\item Perform periodic social engineering tests and communicate summary results
			\item Implement a DLP system that points users to corporate policies
		\end{enumerate}
	\end{frame}
	\begin{frame}{D4-Q194: Answer}
		\begin{block}{Correct Answer: C}
			Users who are aware of security threats may need a reminder that these threats are real. Periodic social engineering tests help maintain a level of alertness. Disciplinary processes are inappropriate for spear phishing.
		\end{block}
	\end{frame}
	
	% ---- D4 Q195 ----
	\begin{frame}{D4-Q195: Increased Malware Attacks Recommendation}
		\begin{block}{Question}
			IT operations indicates a heavy increase in malware attacks. Which is the MOST appropriate recommendation?
		\end{block}
		\begin{enumerate}[(A)]
			\item Contract with a new anti-malware vendor
			\item Close down the Internet connection to prevent employees from visiting infected sites
			\item Make the number of malware attacks part of each employee's performance metrics
			\item Increase employee awareness training, including roles and responsibilities
		\end{enumerate}
	\end{frame}
	\begin{frame}{D4-Q195: Answer}
		\begin{block}{Correct Answer: D}
			Employee awareness training helps the enterprise avoid and more quickly detect malware attacks, particularly when staff understand typical symptoms and know the incident reporting process.
		\end{block}
	\end{frame}
	
	% ---- D4 Q196 ----
	\begin{frame}{D4-Q196: When Data Warehouse Should Remain In-House}
		\begin{block}{Question}
			Which MOST likely indicates that a customer data warehouse should remain in-house rather than outsourced offshore?
		\end{block}
		\begin{enumerate}[(A)]
			\item Telecommunications costs may be higher in the first year
			\item Privacy laws may prevent a cross-border flow of information
			\item Time zone differences may impede communications
			\item Software development may require more detailed specifications
		\end{enumerate}
	\end{frame}
	\begin{frame}{D4-Q196: Answer}
		\begin{block}{Correct Answer: B}
			Privacy laws prohibiting the cross-border flow of personally identifiable information may make it impossible to locate a data warehouse containing customer information in another country.
		\end{block}
	\end{frame}
	
	% ---- D4 Q197 ----
	\begin{frame}{D4-Q197: Responsible for Information Classification}
		\begin{block}{Question}
			The person responsible for ensuring that information is classified is the:
		\end{block}
		\begin{enumerate}[(A)]
			\item Security manager
			\item Technology group
			\item Data owner
			\item Senior management
		\end{enumerate}
	\end{frame}
	\begin{frame}{D4-Q197: Answer}
		\begin{block}{Correct Answer: C --- Data owner}
			The data owner is responsible for applying the proper classification to the data. The security manager applies protection; the technology group may have custody; senior management is responsible for the enterprise.
		\end{block}
	\end{frame}
	
	% ---- D4 Q198 ----
	\begin{frame}{D4-Q198: Transmitting Personal Information}
		\begin{block}{Question}
			Which is the MOST important when transmitting personal information across networks?
		\end{block}
		\begin{enumerate}[(A)]
			\item Encrypting the personal information
			\item Obtaining consent to transfer personal information
			\item Ensuring the privacy of the personal information
			\item Ensuring effective change management
		\end{enumerate}
	\end{frame}
	\begin{frame}{D4-Q198: Answer}
		\begin{block}{Correct Answer: C}
			Privacy protection is necessary to ensure that the receiving party has the appropriate level of protection of personal data. Encryption is a partial answer; consent is not always required; change management is supportive.
		\end{block}
	\end{frame}
	
	% ---- D4 Q199 ----
	\begin{frame}{D4-Q199: Excessive Workstation Images KRI}
		\begin{block}{Question}
			An excessive number of standard workstation images can be categorized as a key risk indicator for:
		\end{block}
		\begin{enumerate}[(A)]
			\item Change management
			\item Configuration management
			\item IT operations management
			\item Data management
		\end{enumerate}
	\end{frame}
	\begin{frame}{D4-Q199: Answer}
		\begin{block}{Correct Answer: B --- Configuration management}
			An excessive number of unique workstation images indicates poor configuration management processes and insufficient attention to actual business requirements during image definition.
		\end{block}
	\end{frame}
	
	% ---- D4 Q200 ----
	\begin{frame}{D4-Q200: Addressing Data Leakage}
		\begin{block}{Question}
			Which BEST addresses the risk of data leakage?
		\end{block}
		\begin{enumerate}[(A)]
			\item Capacity planning procedures
			\item File backup procedures
			\item Acceptable use policies
			\item Database integrity checks
		\end{enumerate}
	\end{frame}
	\begin{frame}{D4-Q200: Answer}
		\begin{block}{Correct Answer: C --- Acceptable use policies}
			Acceptable use policies are an important control for preventing unauthorized disclosure of confidential information. Capacity planning and integrity checks do not address confidentiality; backup procedures are not the best data leakage control.
		\end{block}
	\end{frame}
	
	% ============================================================
	% FINAL SLIDES
	% ============================================================
	\section{Final Review}
	
	\begin{frame}{Exam Summary: 200 Hard Questions}
		\begin{block}{Domain Coverage}
			\begin{itemize}
				\item \textbf{Domain 1: Governance} --- Q1--Q55 (55 questions)
				\item \textbf{Domain 2: IT Risk Assessment} --- Q56--Q100 (45 questions)
				\item \textbf{Domain 3: Risk Response and Reporting} --- Q101--Q165 (65 questions)
				\item \textbf{Domain 4: IT and Security} --- Q166--Q200 (35 questions)
			\end{itemize}
		\end{block}
		\begin{exampleblock}{How to Use This Deck}
			Read each question, commit to an answer, then advance to the next slide to verify against the ISACA justification.
		\end{exampleblock}
	\end{frame}
	
	\begin{frame}{Key Takeaways Across Domains}
		\begin{block}{Golden Rules from the QAE Manual}
			\begin{itemize}
				\item Business management owns risk; senior management accepts it; the board is ultimately accountable
				\item Risk identification precedes assessment; you cannot protect what you do not know exists
				\item Residual risk = remaining risk after controls; acceptable risk = residual within tolerance
				\item BIA establishes RTO, RPO, MTD; risk assessment evaluates likelihood and impact
				\item Control types: preventive (stop), detective (find), corrective (fix), deterrent (discourage), compensating (alternate)
				\item KRIs predict risk; KPIs measure performance; KCIs measure control effectiveness
				\item Cost-benefit analysis justifies controls; when cost exceeds benefit, accept the risk
				\item Encryption = confidentiality; hashing = integrity; digital signature = integrity + nonrepudiation
			\end{itemize}
		\end{block}
	\end{frame}
	
	\begin{frame}{Final Advice}
		\begin{block}{Exam Strategy}
			\centering
			\vspace{0.3cm}
			\begin{itemize}
				\item Read the entire stem; look for \textbf{BEST}, \textbf{MOST}, \textbf{FIRST}
				\item Eliminate wrong choices, then choose the \textbf{business-focused} answer
				\item Understand the \textit{why} behind each justification
				\item Practice scenario-based reasoning --- application beats memorization
			\end{itemize}
			\vspace{0.5cm}
			\Large{\textbf{Good Luck on Your CRISC Exam!}}
		\end{block}
	\end{frame}
	
\end{document}